\chapter*{等式与不等式专题练习}
\begin{enumerate}
    \item 根与系数的关系
    \begin{enumerate}
        \item 已知方程$2x^2+4x-3=0$的两个根为$x_1, x_2$，求下列各式的值：
        \begin{enumerate}
            \item $x_1^2x_2 + x_2^2x_1$
            \item $\dfrac{1}{x_1} + \dfrac{1}{x_2}$
            \item $x_1^2 + x_2^2$
            \item $x_1^3 + x_2^3$
            \item $|x_1 - x_2|$
        \end{enumerate}
        \begin{jieda}
            \begin{enumerate}
                \item $x_1^2x_2 + x_2^2x_1 = x_1x_2(x_1 + x_2) = \dfrac{-3}{2} \times (-2) = 
                3$
                \item $\dfrac{1}{x_1} + \dfrac{1}{x_2} = \dfrac{x_1+x_2}{x_1x_2} = \dfrac{-2}{\dfrac{-3}{2}} = \dfrac{4}{3}$
                \item $x_1^2 + x_2^2 = (x_1+x_2)^2 - 2x_1x_2 = 4 - (-3) = 7$
                \item $x_1^3 + x_2^3 = (x_1 + x_2)(x_1^2 + x_2^2 - x_1x_2) = (-2)(7 - (-\dfrac{3}{2})) = -17$
                \item $\sqrt{(x_1 - x_2)^2} = \sqrt{(x_1 + x_2)^2 - 4x_1x_2} = \sqrt{(-2)^2 - (-6)} = \sqrt{10}$
            \end{enumerate}
            提示：$a^3+b^3 = (a+b)(a^2+b^2-ab), \quad a^3 - b^3 = (a-b)(a^2+b^2+ab)$
        \end{jieda}
        \item 若关于$x$的方程$x^2-(2m-1)x+m^2-2 = 0$的两个实数根分别是$x_1, x_2$，且$x_1^2 - x_1x_2 + x_2^2 = 12$，那么$m = $\blkda{-1}
        \begin{jieda}
            $x_1^2 - x_1x_2 + x_2^2 = 12$即$(x_1+x_2)^2 - 3x_1x_2 = (2m-1)^2 -3(m^2-2) = 12$，解得$m = -1$或5。\\
            验证原方程判别式$\Delta = (2m-1)^2 - 4(m^2-2) \geq 0$，得到$m \leq \dfrac{9}{4}$ \\
            因此有$m = -1$
        \end{jieda}
        \item 设$x_1, x_2$是方程$x^2+x-3=0$的两个实数根，则$x_1^2-x_2+2021 = $\blkda{$2025$}
        \begin{jieda}
            方程的根可以使得方程的等式成立，故$x_1^2+x_1-3=0$，即$x_1^2 = 3 - x_1$，故$x_1^2-x_2+2021 = 2024 - (x_1 + x_2) = 2025$
        \end{jieda}
        \item $\alpha, \beta$是方程$4x^2-4kx+3k+10 = 0$的两个实根，则$(\alpha - 1)^2 + (\beta - 1)^2$的最小值是\blkda{$\dfrac{9}{2}$}
        \begin{jieda}
            $\Delta = 16k^2 - 16(3k+10) \geq 0$，即$k^2 - 3k - 10 \geq 0$，解得$k \in (-\infty, -2] \cup [5, +\infty)$，根据韦达定理有$\left \{ \begin{array}{l}
                 \alpha + \beta = k  \\
                 \alpha \cdot \beta = \dfrac{3k+10}{4}
            \end{array}\right .$ \\
            $(\alpha - 1)^2 + (\beta - 1)^2 = \alpha^2 + \beta^2 + 2 - 2(\alpha + \beta) = (\alpha + \beta)^2 - 2(\alpha + \beta) + 2 - 2\alpha \cdot \beta = k^2 - \dfrac{7k}{2} - 3$ \\
            设$f(k) = k^2 - \dfrac{7k}{2} - 3$，对称轴为$k = \dfrac{7}{4}$，故最小值为$f(5) = \dfrac{9}{2}$.
        \end{jieda}
        \item 若$f(x) = (x^2-10x+c_1)(x^2-10x+c_2)(x^2-10x+c_3)(x^2-10x+c_4)(x^2-10x+c_5)$，且$c_1 \geq c_2 \geq c_3 \geq c_4 \geq c_5 > 0$. 若集合$M = \{x | f(x) = 0\} = \{x_1, x_2, x_3, ...,x_9\} \subseteq \mathbf{N}$，则$c_1 - c_5 = $\blkda{16}
        \begin{jieda}
            因为$M$有9个元素，因此$x^2-10x+c_i = 0$中有一个有重根，又因为二次项系数为正，所以必为$x^2-10x+c_1 = 0$有重根$x = 5$，$\Delta = 100 - 4c_1 = 0$，故$c_1 = 25$. 根据韦达定理可知$M$中其他元素关于$5$成对出现，又因为$M \subseteq \mathbf{N}$，故$M = \{1,2,3,4,5,6,7,8,9\}$，所以$x^2-10x+c_5 = 0$的解集为$\{1,9\}$.根据韦达定理有$c_5 = 9$，因此$c_1 - c_5 = 16$.
        \end{jieda}
        \item 已知关于$x$的方程$x^2-(2k+1)x+4\left(k-\dfrac{1}{2}\right)= 0$.
        \begin{enumerate}
            \item 求证：无论$k$取何值，此方程总有实数根.
            \item 若等腰三角形$ABC$的一边长$a = 4$，另两边的长$b, c$恰是这个方程的两根，求三角形$ABC$的周长.
        \end{enumerate}
        \begin{jieda}
            \begin{enumerate}
                \item $\Delta = (2k+1)^2 - 16(k-\dfrac{1}{2}) = 4k^2 + 4k + 1 - 16k + 8 = 4k^2 - 12k + 9$，$\Delta_k = 0$，故$\Delta \geq 0$恒成立，因此无论$k$取何值，此方程总有实数根.
                \item \ding{172} $b = a = 4$，根据韦达定理有$c = k-\dfrac{1}{2}$，根据韦达定理进一步有$4+k-\dfrac{1}{2} = 2k+1$，解得$k = \dfrac{5}{2}$，故$c = 2$，满足三角形要求，此时周长为10. \\
                \ding{173} $c = a = 4$和上一种情况对称，此时周长为10. \\
                \ding{174} $b = c$，此时有$\Delta = 0$，即$k = \dfrac{3}{2}$，故$b = c = 2$，此时不满足三角形要求. \\
                综上，周长为10
            \end{enumerate}
        \end{jieda}
    \end{enumerate}
    \item 从直线位置关系看二元一次方程组
    \begin{enumerate}
        \item 已知关于$x, y$的方程组$\left \{ \begin{array}{l}
         ax+2y = 1+a  \\
         2x+2(a-1)y = 3 
        \end{array} \right .$，当$a$为何值时，方程有唯一解，无解，无穷多组解？
        \begin{jieda}
            \begin{dingautolist}{172}
                \item 当$a = 1$时，方程组化为$\left \{ \begin{array}{l}
                     y = -\dfrac{x}{2} + 1  \\
                     x = \dfrac{3}{2}
                \end{array} \right .$ 此时有唯一解
                \item 当$a \neq 1$时，将方程都化为斜截式：
                $\left \{ \begin{array}{l}
                     y = -\dfrac{ax}{2} + \dfrac{1+a}{2}  \\
                     y = \dfrac{-2x}{2(a-1)} + \dfrac{3}{2(a-1)}
                \end{array} \right .$ ($a \neq 1$) \\
                当斜率不等时，方程有唯一解。即$a \neq -1$ 或 2时方程有唯一解$\left \{ \begin{array}{l}
                     x = \dfrac{a+2}{a+1}  \\
                     y = \dfrac{1}{2(a+1)}
                \end{array} \right .$
            \end{dingautolist}
            综上，当$a \neq -1$ 或 2时方程有唯一解
            
            当斜率相等时，即$a = -1$ 或 2，进一步要判断截距是否相等
            \begin{dingautolist}{172}
                \item 若$a = -1$时，方程组无解
                \item 若$a = 2$时，方程组有无穷多解
            \end{dingautolist}
        \end{jieda}
        \item 已知方程$|x| = ax+1$有一负解，且无正解，则实数$a$的取值范围为\blkda{$[1, +\infty)$}
        \begin{jieda}
            $y = |x|$对应着直线$y = x$与$y = -x$在$x$轴及其以上的部分。而$y = ax+1$对应着一条过$(0, 1)$的直线，随着斜率$a$的变化，直线的位置关系发生着改变。\\
            方程$|x| = ax+1$有一负解，且无正解则意味着两者只在$y$轴左侧有交点，因此斜率$a \geq 1$
        \end{jieda}
        \item 已知集合$A = \left\{(x, y) | \dfrac{y - 3}{x - 2} = a + 1 \right \}, B = \left\{(x, y) | (a^2 - 1)x + (a-1)y = 15\right\}$，当$a$为何实数的时，$A \cap B = \varnothing$
        \begin{jieda}
            从形式上看$A, B$两个集合都是点集。\\
            集合$A$对应于直线$y = (a+1)(x-2) + 3$，但不包括其上的点$(2, 3)$ \\
            集合$B$对应于直线$y = -(a+1)x + \dfrac{15}{a-1} (a \neq 1)$，而$B = \varnothing (a = 1)$ 
            \begin{enumerate}
                \item 当$a = 1$时，由于$B = \varnothing$，自然满足条件
                \item 当$a \neq 1$时，若两直线平行或直线$y = -(a+1)x + \dfrac{15}{a-1}$过$(2, 3)$则满足条件
                \begin{dingautolist}{172}
                    \item 当$a = -1$时，两直线平行
                    \item 当$a = -4 \mbox{或} \dfrac{5}{2}$的时候直线$y = -(a+1)x + \dfrac{15}{a-1}$过$(2, 3)$
                \end{dingautolist}
            \end{enumerate}
            综上$a \in \left\{1, -1, -4, \dfrac{5}{2}\right\}$
        \end{jieda} 
    \end{enumerate}
    \item 直线过定点问题
    \begin{enumerate}
        \item 对于任意实数$k$，$y = kx+(2k-2)$恒经过的定点是？
        \begin{jieda}
            既然要对任意实数k都存在这样的性质，就应该让k的系数为0才可以使得k的变化不影响该性质。因此若将k看做主元，得到关于k的一元一次方程后让k的系数和常数项均为0，形成0=0的恒等式。 \\
            $y = kx+(2k-2)$即$(x+2)k-(2+y) = 0$，恒过$(-2, -2)$
        \end{jieda}
        \item 对于任意实数$k$，$(k+1)x + (1-k)y + (5-k) = 0$恒经过的定点是？
        \begin{jieda}
            道理同上一题，$(k+1)x + (1-k)y + (5-k) = 0$即$(x-y-1)k+(x+y+5) = 0$，恒过$(-2, -3)$
        \end{jieda}
    \end{enumerate}
    \item 根据不等式的性质比较大小
    \begin{enumerate}
        \item 设$ab>0$，且$\dfrac{c}{a} > \dfrac{d}{b}$，则下列各式中，恒成立的是\dotfill{(\kh{B})}
        \xx{$bc < ad$}{$bc > ad$}{$\dfrac{a}{c} > \dfrac{b}{d}$}{$\dfrac{a}{c} < \dfrac{b}{d}$}
        \begin{jieda}
            $\dfrac{c}{a} > \dfrac{d}{b} \Longrightarrow \dfrac{c}{a} \cdot (ab) > \dfrac{d}{b}  \cdot (ab) \Longrightarrow ac > bd$
        \end{jieda}
        \item 下列命题中，不正确的一个是\dotfill{(\kh{D})}
        \xx{若$\sqrt[3]{a} > \sqrt[3]{b}$，则$a>b$}{若$a > b, c > d$，则$a - d > b - c$}{若$a > b > 0, c>d>0$，则$\dfrac{a}{d} > \dfrac{b}{c}$}{若$a>b>0, ac > bd$，则$c>d$}
        \begin{jieda}
            \begin{enumerate}[A.]
                \item 若$a > 0, b \leq 0$，或$a \geq 0, b < 0$显然成立\\
                若$ab > 0$，则$a - b = (\sqrt[3]{a} - \sqrt[3]{b})((\sqrt[3]{a})^2 + (\sqrt[3]{b})^2 + \sqrt[3]{a} \cdot \sqrt[3]{b}) > 0$ 
                
                （换元更好）
                \item $c > d \Longrightarrow -d > -c$，故$a - d > b - c$
                \item $a > b > 0, c>d>0 \Longrightarrow ac - bd > 0 \Longrightarrow \dfrac{ac-bd}{cd} > 0 \Longrightarrow \dfrac{a}{d} - \dfrac{b}{c} > 0$
                \item 反例：$a = 10, b = 1, c = \dfrac{1}{2}, d = 1$
            \end{enumerate}
        \end{jieda}
        \item 若$x < y < 0$，则有\dotfill{(\kh{B})}
        \xx{$0 < x^2 < xy$}{$y^2 < xy < x^2$}{$xy < y^2 < x^2$}{$0 < x^2 < y^2$}
        \begin{jieda}
            $x < y < 0 \Longrightarrow |x| > |y|$
        \end{jieda}
        \item 用不等号（“$>$”或“$<$”）填空：
        \begin{enumerate}
            \item 若$a \neq b$，则$a^2+3b^2$ \blkda{>} $2b(a+b)$
            \item 若$c > 1$，则$\sqrt{c+1} - \sqrt{c}$ \blkda{<} $\sqrt{c} - \sqrt{c-1}$
            \item 若$a > b, c > d$，且$a, d$都是负数，则$ac$ \blkda{<} $bd$
            \item 若$a > 0, b > 0$且$a \neq b$，则$\dfrac{a^2}{b} + \dfrac{b^2}{a}$\blkda{>} $a+b$
        \end{enumerate}
        \begin{jieda}
            \begin{enumerate}
                \item $LHS - RHS = a^2 - 2ab + b^2 = (a-b)^2 > 0$
                \item $LHS - RHS = \dfrac{1}{\sqrt{c+1} + \sqrt{c}} - \dfrac{1}{\sqrt{c} + \sqrt{c-1}} = \dfrac{\sqrt{c-1} - \sqrt{c+1}}{(\sqrt{c+1} + \sqrt{c})(\sqrt{c} + \sqrt{c-1})} < 0$
                \item 若$c \geq 0$，则$LHS \leq 0 < RHS$，若$c < 0$，则$LHS = (-a)(-c) < (-b)(-d) = RHS$
                \item $LHS - RHS = \dfrac{a^3+b^3}{ab} - (a+b) = \dfrac{(a+b)}{ab}[(a^2+b^2-ab) - ab]  = \dfrac{(a+b)(a-b)^2}{ab} > 0$
            \end{enumerate}
        \end{jieda}
        \item 已知$a \in \mathbf{R}$，比较$\dfrac{1}{a+1}$与$1-a$的大小.
        \begin{jieda}
            $\dfrac{1}{a+1} - (1-a) = \dfrac{a^2}{a+1}$，故\ding{172}$a < -1$时，$\dfrac{1}{a+1} < 1-a$；\ding{173}$a = 0$时，$\dfrac{1}{a+1} = 1-a$；\ding{174}$a \in (-1, 0) \cup (0, +\infty)$时，$\dfrac{1}{a+1} > 1-a$；
        \end{jieda}
    \end{enumerate}
    \item 根据同向可加性，同向同正可乘性求取值范围
    \begin{enumerate}
        \item 已知$-1 < x-y < 4, 2 < x+y < 3$，则$3x+y$的取值范围是\blkda{$(3, 10)$}
        \begin{jieda}
            设$3x+y = a(x-y) + b(x+y)$，解得$a = 1, b = 2$ \\
            $2(x+y) \in (4, 6)$，因此$3x+y \in (3, 10)$
        \end{jieda}
        \item 已知$30 < x < 42, 16 < y < 24$，求$x+y, x-2y, \dfrac{x}{y}$的取值范围.
        \begin{jieda}
            \begin{dingautolist}{172}
                \item $46<x+y<66$
                \item $16 < y < 24 \Longrightarrow -48 < -2y < -32$，故$-18 < x-2y < 10$
                \item $16 < y < 24 \Longrightarrow \dfrac{1}{24} < \dfrac{1}{y} < \dfrac{1}{16}$，故$\dfrac{5}{4} < \dfrac{x}{y} < \dfrac{21}{8}$
            \end{dingautolist}
        \end{jieda}
        \item 已知$-1 \leq a+b \leq 1$，$1 \leq a-b \leq 3$，$10a - b$的取值范围是\blkda{$[1,21]$}.
        \begin{jieda}
            设$10a - b = x(a+b) + y(a-b)$，即$\left \{ \begin{array}{l}
                x+y = 10  \\
                x-y = -1 
            \end{array} \right.$，解得$x = \dfrac{9}{2}, y = \dfrac{11}{2}$，故$-\dfrac{9}{2} \leq \dfrac{9}{2}(a+b) \leq \dfrac{9}{2}$，$\dfrac{11}{2} \leq \dfrac{11}{2}(a-b) \leq \dfrac{33}{2}$，因此$1 \leq 10a-b \leq 21$.
        \end{jieda}
        \item 已知实数$x, y$满足$3 \leq xy^2 \leq 8, 4 \leq \dfrac{x^2}{y} \leq 9$，$\dfrac{x^3}{y^4}$的取值范围是\blkda{$[2, 27]$}.
        \begin{jieda}
            设$\dfrac{x^3}{y^4} = (xy^2)^a \cdot (\dfrac{x^2}{y})^b$，即$\left \{ \begin{array}{l}
                a+2b = 3  \\
                2a-b = -4 
            \end{array} \right.$，解得$a = -1, b = 2$，故$\dfrac{1}{8} \leq \dfrac{1}{xy^2} \leq \dfrac{1}{3}$，$16 \leq \dfrac{x^4}{y^2} \leq 81$，因此$2 \leq \dfrac{x^3}{y^4} \leq 27$.
        \end{jieda}
    \end{enumerate}
    \item 求解高次不等式、分式不等式
    \begin{enumerate}
        \item 解不等式：
        \begin{dingautolist}{172}
            \item $(x-3)(x+1) > 0$
            \item $(x-1)(x+4) < 0$
            \item $-x^2+3x-\dfrac{1}{2} \geq 0$
            \item $x^2 - 5x + 6 > 0$
            \item $9x^2 - 6x + 1 > 0$
            \item $-x^2 + 2x - 3 > 0$
            \item $x^2 - x + \dfrac{1}{4} > 0$
            \item $x^2 - 3x +4 > 0$
        \end{dingautolist}
        \begin{jieda}
            \begin{dingautolist}{172}
                \item $(-\infty, -1) \cup (3, +\infty)$
                \item $(-4, 1)$
                \item $\left[\dfrac{3-\sqrt{7}}{2}, \dfrac{3+\sqrt{7}}{2}\right]$
                \item $(-\infty, 2) \cup (3, +\infty)$
                \item $\left\{x|x \neq \dfrac{1}{3}\right\}$
                \item $\varnothing$
                \item $\left\{x|x \neq \dfrac{1}{2}\right\}$
                \item $\mathbf{R}$
            \end{dingautolist}
        \end{jieda}
        \item 设命题$P$：关于$x$的不等式$a_1x^2+b_1x+c_1 > 0$与不等式$a_2x^2+b_2x+c_2 > 0$的解集相同；命题$Q$：$\dfrac{a_1}{a_2} = \dfrac{b_1}{b_2} = \dfrac{c_1}{c_2}$。则命题$Q$是命题$P$的\dotfill{(\kh{D})}.
        \xx{充要条件}{充分不必要条件}{必要不充分条件}{既不充分也不必要条件}
        \begin{jieda}
            如果两个不等式的解集都是实数集或者空集，则系数不一定成比例。如果系数比例为-1，则解集“互补”（忽略端点开闭）。因此选D。
        \end{jieda}
        \item 解不等式
        \begin{dingautolist}{172}
            \item $\dfrac{x+3}{4-x} > 0$
            \item $\dfrac{5x+3}{x-1} \leq 3$
            \item $\dfrac{x+5}{x^2+2x+3} \leq 1$
            \item $\dfrac{3-2x}{x-1} < 0$
            \item $\dfrac{4+x}{2+x} \geq 2$
            \item $\dfrac{4-x}{x^2+x+1} \leq -1$
        \end{dingautolist}
        \begin{jieda}        
            \begin{dingautolist}{172}
                \item $x \in (-3, 4)$
                \item $x \in [-3, 1)$
                \item $x \in (-\infty, -2] \cup [1, +\infty)$
                \item $x \in (-\infty, 1) \cup (\dfrac{3}{2}, +\infty)$
                \item $x \in (-2, 0]$
                \item 解集为$\varnothing$
            \end{dingautolist}
        \end{jieda}
        \item 不等式$\dfrac{x+5}{x^2+2x+3} \geq 1$的解集为\blkda{$[-2, 1]$}
        \begin{jieda}
            因为$x^2+2x+3 > 0$恒成立，因此可以转化为$x+5 \geq x^2+2x+3$，解集为$[-2, 1]$ \\
            \ps 先判断分母是否恒正或者恒负，再考虑移项通分.
        \end{jieda}
        \item 不等式$\dfrac{x^2(x+3)(x^2-x+1)}{x^2 - 5x + 4} > 0$的解集为\blkda[5]{$(-3, 0) \cup (0, 1) \cup (4, +\infty)$}
        \begin{jieda}
            \ps 穿针引线法求解，奇穿偶回，\textbf{偶次方根绝对不能丢弃，且易出错}.
        \end{jieda}
        \item 已知集合$A = \{x | \dfrac{2x-1}{x^2+3x+2} > 0\}, B = \{x | x^2 + ax + b \leq 0\}$，且$A \cap B = \{x | \dfrac{1}{2} < x \leq 3\}$，求实数$a, b$的取值范围.
        \begin{jieda}
            解分式不等式有$A = (-2, -1) \cup (\dfrac{1}{2}, +\infty)$，设$B = [m, n]$，则根据$A \cap B = \{x | \dfrac{1}{2} < x \leq 3\}$可知$n = 3, m \in [-1, \dfrac{1}{2}]$，根据韦达定理有$a = -(3+m) \in \left[-\dfrac{7}{2}, -2\right]$，$b = 3m \in \left[-3, \dfrac{3}{2}\right]$
        \end{jieda}
        \item 已知$a, b \in \mathbf{R}$且$ab \neq 0$，对任意$x \geq 0$均有$(x-a)(x-b)(x-2a-b) \geq 0$，则\dotfill{(\kh{C})}
        \xx{$a<0$}{$a>0$}{$b < 0$}{$b > 0$}
        \begin{jieda}
            穿针引线绘制出函数图像草图。函数图像和$x$轴有三个公共点，横坐标为$a, b, 2a+b$。
            \begin{enumerate}
                \item 若三个点均在原点左侧，则满足题意，此时$a < 0, b < 0$
                \item 若有点在$x$轴右侧，则必然左侧有一个点，右侧有两个重合的点才能满足题意，此时绘制函数图像草图可知函数图像和$x$轴在两个重合的点出相切。唯一的可能是$a > 0, b = -a$，此时$a$和$2a+b$相等。
            \end{enumerate}
            综上必有$b < 0$
        \end{jieda}
    \end{enumerate}
    \item 最高次含参导致的退化问题
    \begin{enumerate}
        \item 集合$A = \{x | x^2-8x+12 = 0\}, B = \{x | ax-1 = 0\}$，且$A \cup B = A$，则$a$的集合为\blkda{$\left\{0, \dfrac{1}{2}, \dfrac{1}{6}\right\}$}
        \begin{jieda}
            $A \cup B = A$即$B \subseteq A$
            \begin{dingautolist}{172}
                \item $a = 0$，此时$B = \varnothing$，满足要求.
                \item $a \neq 0$，此时$B = \left \{\dfrac{1}{a}\right\}$，因为$A = \{2, 6\}$，若$B \subseteq A$，则$\dfrac{1}{a} \in A$，因此$a = \dfrac{1}{2} \mbox{或} \dfrac{1}{6}$
            \end{dingautolist}
        \end{jieda}
        
        \item 解关于$x$的不等式$ax^2-(a-2)x-2 > 0$
        \begin{jieda}
            \begin{dingautolist}{172}
                \item $a = 0$，不等式化为$2x-2 > 0$，解集为$(1, +\infty)$
                \item $a < 0$，对应的方程$ax^2-(a-2)x-2 = 0$的两根为$x_1 = 1, x_2 = -\dfrac{2}{a}$ \\
                \begin{enumerate}[1.]
                    \item $0 > a > -2$时解集为$\left(1, -\dfrac{2}{a}\right) $
                    \item $a = -2$时解集为$\varnothing$
                    \item $a < -2$时解集为$\left(-\dfrac{2}{a}, 1\right) $
                \end{enumerate}
                \item $a > 0$，对应的方程$ax^2-(a-2)x-2 = 0$的两根为$x_1 = 1, x_2 = -\dfrac{2}{a}$，因此解集为$\left(-\infty, -\dfrac{2}{a}\right) \cup (1, +\infty) $
            \end{dingautolist}
        \end{jieda}
        \item 对任意实数$x$，不等式$2kx^2+kx-3<0$恒成立，则实数$k$的取值范围是\blkda{$(-24, 0]$}
        \begin{jieda}
            当$k=0$时，不等式成立。\\
            当$k \neq 0$，从二次函数的角度看一元二次不等式，要使得不等式恒成立，则应有函数开口向下，即$k<0$，并且有$\Delta = k^2 + 24k < 0$，因此$k \in (-24, 0)$。\\
            综上，$k \in (-24, 0]$
        \end{jieda}
         \item 若关于$x$的方程$kx^2+2x-1 = 0$有两个不相等的实数根，则$k$的取值范围是\blkda[3]{$(-1, 0) \cup (0, +\infty)$}
        \begin{jieda}
            若$k = 0$，则方程退化为一次方程，不会有两根。\\
            若$k \neq 0$，则应有$\Delta = 4+4k > 0$，即$k > -1$，故$k \in (-1, 0) \cup (0, +\infty)$
        \end{jieda}
        \item 若关于$x$的方程$2kx^2+(8k+1)x+8k = 0$有两个不等实根，则实数$k$的取值范围是\blkda[3.5]{$\left(-\dfrac{1}{16}, 0\right) \cup (0, +\infty)$}
        \begin{jieda}
            若$k = 0$，则方程退化为一次方程，不会有两根。\\
            若$k \neq 0$，则应有$\Delta = (8k+1)^2 - 64k^2 = 16k + 1 > 0$，因此$k \in \left(-\dfrac{1}{16}, 0\right) \cup (0, +\infty)$
        \end{jieda}
    \end{enumerate}
    \item 根据解集求参数
    \begin{enumerate}
        \item 若关于$x$的不等式$ax^2+bx-2>0$的解集是$\left(-\infty, -\dfrac{1}{2}\right) \cup \left(\dfrac{1}{3}, +\infty\right)$，则$ab = $\blkda{24}
        \begin{jieda}
            从解集的形式可知\ding{172} $a > 0$，\ding{173}$ax^2+bx-2=0$的两根是$-\dfrac{1}{2}$和$\dfrac{1}{3}$，将两根代入解得$a = 12, b = 2$，因此$ab = 24$
        \end{jieda}
        \item 若不等式$(2a-b)x +3a - 4b < 0$的解集是$\left(\dfrac{9}{4}, +\infty\right)$，则不等式$(a-4b)x + 2a - 3b > 0$的解集是\blkda[2]{$\left(-\dfrac{8}{19}, +\infty\right)$}
        \begin{jieda}
            从解集形式可知\ding{172} $2a-b < 0$，\ding{173}$(2a-b)x +3a - 4b = 0$的根是$\dfrac{9}{4}$，即$a = \dfrac{5}{6}b$，故$(a-4b)x + 2a - 3b > 0$可转化为$\left(-\dfrac{19}{5}a\right)x - \left(\dfrac{8}{5}a\right) > 0$，因为$2a-b < 0$，所以$a < 0$，故可以进一步转化为$\dfrac{19}{5}x + \dfrac{8}{5} < 0$，解得$x \in \left(-\dfrac{8}{19}, +\infty\right)$
        \end{jieda}
       
        \item 若关于$x$的不等式$ax^2-2ax+2a+3>0 (a \neq 0)$无实数解，则$a$的取值范围是\blkda{$(-\infty, -3]$}
        \begin{jieda}
            $a \neq 0$，应有$a < 0$，否则对应的二次函数开口向上必然有实数解. 此时应有$\Delta = 4a^2 - 4a(2a+3) = -4a^2 - 12a \leq 0$，解得$a \in (-\infty, -3] \cup [0, +\infty)$. 综上，$a \in (-\infty, -3]$
        \end{jieda}
        \item 若关于$x$的不等式$2x-1 > a(x-2)$的解集是$\mathbf{R}$，则实数$a$的取值范围是\blkda{$\{2\}$}
        \begin{jieda}
            问题转化为直线$y = 2x-1$恒在直线$y = a(x-2)$上方，当且仅当两直线平行时满足要求，即$a = 2$.
        \end{jieda}
    \end{enumerate}
    \item 同解不等式
    \begin{enumerate}
        \item 若关于$x$的不等式$ax^2+bx+c > 0$的解集是$\{x|3<x<5\}$，则不等式$cx^2+bx+a<0$的解集是\blkda[3.5]{$\left(-\infty, \dfrac{1}{5}\right) \cup \left(\dfrac{1}{3}, +\infty\right)$}
        \begin{jieda}
            因为$0$不在解集中，故$a+\dfrac{b}{x}+\dfrac{c}{x^2} > 0$的解集是$\{x|3<x<5\}$，令$t = \dfrac{1}{x}$，则$t \in \left(\dfrac{1}{5}, \dfrac{1}{3}\right)$时$x \in (3, 5)$，使得$a+\dfrac{b}{x}+\dfrac{c}{x^2} > 0$成立，即$a+bt+ct^2 > 0$成立，故$cx^2+bx+a>0$解集为$\left(\dfrac{1}{5}, \dfrac{1}{3}\right)$，即$cx^2+bx+a<0$的解集是$\left(-\infty, \dfrac{1}{5}\right) \cup \left(\dfrac{1}{3}, +\infty\right)$
        \end{jieda}
        \item 已知关于$x$的不等式$\dfrac{k}{x+a} + \dfrac{x+b}{x+c} < 0$的解集为$(-2, -1) \cup (2, 3)$，关于$x$的不等式$\dfrac{kx}{ax-1} + \dfrac{bx-1}{cx-1} < 0$的解集为\blkda[3.5]{$\left(-\dfrac{1}{2}, -\dfrac{1}{3}\right) \cup \left(\dfrac{1}{2}, 1\right)$}
        \begin{jieda}
            因为$0$不在解集中，故$\dfrac{k}{x+a} + \dfrac{x+b}{x+c} < 0$可以变形为$\dfrac{k\cdot\left(-\dfrac{1}{x}\right)}{(x+a)\cdot\left(-\dfrac{1}{x}\right)} + \dfrac{(x+b)\cdot\left(-\dfrac{1}{x}\right)}{(x+c)\cdot\left(-\dfrac{1}{x}\right)} < 0$，进一步得到原不等式的同解不等式$\dfrac{k\cdot\left(-\dfrac{1}{x}\right)}{a\cdot\left(-\dfrac{1}{x}\right) - 1} + \dfrac{b\cdot\left(-\dfrac{1}{x}\right) - 1}{c\cdot\left(-\dfrac{1}{x}\right) - 1} < 0$ \\
            设$t = -\dfrac{1}{x}$，故当$t \in \left(-\dfrac{1}{3}, -\dfrac{1}{2}\right) \cup \left(\dfrac{1}{2}, 1\right)$，即$-\dfrac{1}{t} \in (-2, -1) \cup (2, 3)$时，满足不等式$\dfrac{kt}{at-1} + \dfrac{bt-1}{ct-1} < 0$。因此关于$x$的不等式$\dfrac{kx}{ax-1} + \dfrac{bx-1}{cx-1} < 0$的解集为$\left(-\dfrac{1}{2}, -\dfrac{1}{3}\right) \cup \left(\dfrac{1}{2}, 1\right)$
        \end{jieda}
    \end{enumerate}
    \item 无理不等式求解
    \begin{enumerate}
        \item 不等式$\sqrt{x^2+x-2} < \dfrac{2x}{3}$的解集为\blkda{$\left[1, \dfrac{6}{5}\right)$}
        \begin{jieda}
            $\sqrt{x^2+x-2} < \dfrac{2x}{3} \Longleftrightarrow \left \{\begin{array}{l}
                 x^2+x-2 \geq 0  \\
                 \dfrac{2x}{3} \geq 0 \\
                 x^2+x-2 < \dfrac{4x^2}{9}
            \end{array} \right.$，解得$x \in \left[1, \dfrac{6}{5}\right)$
        \end{jieda}
        \item 若不等式$2\sqrt{4-x^2} \leq k(x+2)$的解集为$[0, n] \cup \{-2\}$，则$k = $\blkda{2}
        \begin{jieda}
            \begin{enumerate}
                \item 方法1 \\
                $2\sqrt{4-x^2} \leq k(x+2) \Longleftrightarrow \left \{ \begin{array}{l}
                     -2 \leq x \leq 2 \\
                     4(4-x^2) \leq k^2(x+2)^2 \\
                     k(x+2) \geq 0
                \end{array} \right. \\
                \Longleftrightarrow \left \{ \begin{array}{l}
                     -2 \leq x \leq 2 \\
                     (x+2)[(k^2+4)x+2k^2-8] \geq 0 \\
                     k \geq 0
                \end{array} \right. \Longleftrightarrow \left \{ \begin{array}{l}
                     -2 \leq x \leq 2 \\
                     (k^2+4)x+2k^2-8 \geq 0 \\
                     k \geq 0
                \end{array} \right.$或$x = -2$，\\
                即$\left \{ \begin{array}{l}
                     -2 \leq x \leq 2 \\
                     x \geq \dfrac{2(4-k^2)}{k^2+4} \\
                     k \geq 0
                \end{array} \right.$的解集为$[0, n]$，因此必有$\dfrac{2(4-k^2)}{k^2+4} = 0$，即$k = 2$
                \item 方法2 \\
                $2\sqrt{4-x^2} \leq k(x+2)$的解集为$[0, n] \cup \{-2\}$，即$\dfrac{2\sqrt{4-x^2}}{x+2} \leq k$的解集是$[0, n]$，\\
                不等式左侧连续，不等式解集边界是方程的根，带入$x=0$可得$k=2$
            \end{enumerate}
        \end{jieda}
    \end{enumerate}
    \item 求解含绝对值不等式
    \begin{enumerate}
        \item 解不等式：$|x+1| < 3-2x$
        \begin{jieda}
            解集为$\left(-\infty, \dfrac{2}{3}\right)$ \\
            \ps \begin{enumerate*}[1.]
                \item $|f(x)| > g(x)$的充要条件是$f(x) > g(x)$或$f(x) < -g(x)$ 
                \item $|f(x)| < g(x)$的充要条件是$-g(x) < f(x) < g(x)$
            \end{enumerate*} \\
            \ps \textbf{此处绝对不能平方处理}，否则等价于$|x+1| < |3-2x|$会得到$\left(-\infty, \dfrac{2}{3}\right) \cup \left(4, +\infty\right)$；当不等式两侧均含有绝对值时才适合平方去绝对值.
        \end{jieda}
        \item 已知不等式$|x+3| \leq 5$和$(x+6)(x+4)^4(x-1) \geq 0$中有且仅有一个成立，则实数$x$的范围是\blkda[7]{$(-\infty, -8) \cup (-6, -4) \cup (-4, 1) \cup (2, +\infty)$}
        \begin{jieda}
            第一个不等式的解集为$[-8, 2]$，第二个不等式的解集为$(-\infty, -6] \cup [1, +\infty) \cup \{-4\}$ \\
            因此第一个不等式成立第二个不等式不成立时$x \in (-6, -4) \cup (-4, 1)$，第二个不等式成立，第一个不等式不成立时$x \in (-\infty, -8) \cup (2, +\infty)$\\
            两者取并集，得到$(-\infty, -8) \cup (-6, -4) \cup (-4, 1) \cup (2, +\infty)$ 
        \end{jieda}
        \item 关于实数$x$的不等式$\left | x - \dfrac{(a+1)^2}{2} \right | \leq \dfrac{(a-1)^2}{2}$与$|x-a-1| \leq a$的解集依次是$A, B$，求使得$A \subseteq B$的$a$的取值范围.
        \begin{jieda}
            $\left | x - \dfrac{(a+1)^2}{2} \right | \leq \dfrac{(a-1)^2}{2}$即$-\dfrac{(a-1)^2}{2} \leq x - \dfrac{(a+1)^2}{2} \leq \dfrac{(a-1)^2}{2}$，即$2a \leq x \leq a^2+1$ \\
            $|x-a-1| \leq a$ 即 $-a \leq x-a-1 \leq a$，即$1 \leq x \leq 2a+1$ \\
            由于$a^2+1 \geq 2a$恒成立，所以$A$必不为空集，故$A \subseteq B$即$1 \leq 2a, a^2+1 \leq 2a+1$，解得$a \in \left[\dfrac{1}{2}, 2\right]$

        
        
        
        
        
        
        \end{jieda}
        \item 若关于$x$的不等式$ax + 6 + |x^2 - ax - 6| \geq 4$恒成立，$a$的取值范围是\blkda{$[-1, 1]$}
        \begin{jieda}
            原不等式转化为$|x^2 - ax - 6| \geq -2 - ax$恒成立，即$x^2 - ax - 6 \geq -2 - ax$与$x^2 - ax - 6 \leq 2 + ax$解集的并集为$\mathbf{R}$ \\
            前者解集为$(-\infty, -2] \cup [2, +\infty)$，后者解集为$[a-\sqrt{a^2+8}, a + \sqrt{a^2+8}]$，故$(-2, 2) \subseteq [a-\sqrt{a^2+8}, a + \sqrt{a^2+8}]$，因此$a \in [-1, 1]$
        \end{jieda}
        \item 已知函数$f(x) = |x-2| + |x-a|$，若对于任意的$x \in [1, 2]$，$f(x) \geq |x-4|$恒成立，实数$a$的范围是\blkda[4]{$(-\infty, -1] \cup [4, +\infty)$}
        \begin{jieda}
            当$x \in [1, 2]$，$f(x) \geq |x-4|$即$2-x + |x-a| \geq 4-x$，整理得$|x-a| \geq 2$，由绝对值的几何意义可知$a \in (-\infty, -1] \cup [4, +\infty)$ \\
            \ps 有自变量的具体范围的时候，可以先看一下绝对值是否可以直接去掉.
        \end{jieda}
        \end{enumerate}
    \item 包含关系与推出关系结合不等式求解
    \begin{enumerate}
        \item 已知集合$A = \left\{x | \dfrac{x-2}{x+1} < 0\right\}$，$x \in A$的一个必要条件是$x \geq a$，则实数$a$的取值范围是\blkda{$a \leq -1$}
        \begin{jieda}
            $A = (-1, 2)$，故$a \leq -1$ \\
            \ps \textbf{小范围可以推出大范围}
        \end{jieda}
        \item 已知$x \in \mathbf{R}$，“$(x-2)(x-3) \leq 0$”是“$|x-2|+|x-3| = 1$成立”的\blkda{充要}条件.
        \begin{jieda}
            $(x-2)(x-3) \leq 0$解集为$[2, 3]$ \\
            $|x-2|+|x-3| = 1$解集为$[2, 3]$
        \end{jieda}
    \end{enumerate}
    \item 三角不等式
    \begin{enumerate}
        \item 若$ab > 0$，则\ding{172}$|a+b| > |a|$；\ding{173}$|a+b| < |b|$；\ding{174}$|a+b| < |a-b|$;\ding{175}$|a+b| > |a-b|$;\ding{176} $|a-b| < |a| - |b|$；\ding{177} $|a -b| < |a| + |b|$中正确的是\blkda{\ding{172},\ding{175},\ding{177}}
        \begin{jieda}
            $ab > 0$，则$|a+b| = |a| + |b| > |a|$，因此\ding{172}正确，\ding{173}错误. \\
            $ab > 0$，则$|a+b| = |a| + |b| > |a-b|$，因此\ding{175}正确，\ding{174}错误. \\
            \ding{176}和\ding{177}则直接应用三角不等式可以判断.
        \end{jieda}
        \item ``$ab\geq 0$''是``$\abs{a+b}=\abs{a}+\abs{b}$''的\blk 条件. \dotfill(\kh{C})
        \xx
        {充分不必要条件}
        {必要不充分条件}
        {充要条件}
        {既不充分也不必要条件}
        \begin{jieda}
        考察三角不等式的\textbf{取等条件}. \\
        $\abs{a+b} \leq \abs{a}+\abs{b}$的取等条件是$ab\geq 0$ \\
        $\abs{a-b} \leq \abs{a}+\abs{b}$的取等条件是$ab\leq 0$
        \end{jieda}
        \item 对于实数$x,y$，满足$|x+y| \leq 2$且$|x-y| \leq 3$，则$|3x-y|$的最大值为\blkda{8}
        \begin{jieda}
            看题干的形式很容易想到待定系数法，$a(x+y) + b(x-y) = 3x-y$，解得$a = 1, b = 2$，故$|3x-y| = |(x+y) + 2(x-y)| \leq |x+y| + 2|x-y| \leq 8$，当$x = \dfrac{5}{2}, y = -\dfrac{1}{2}$时可以取等。
        \end{jieda}
        \item 已知函数$f(x) = |x+m| + |x-2m|$，若$f(x) \geq 3$恒成立，求$m$的取值范围
        \begin{jieda}
            问题可以转化为$f(x)$的最小值大于等于3恒成立，而 $f(x)$是一个平底锅函数，最小值可以通过三角不等式求得，即 $f(x) = |x+m| + |x-2m| \geq |3m|$，因此有$|3m| \geq 3$，得到$m \in (-\infty, -1] \cup [1, +\infty)$ \\
            \ps 注意\textbf{问题转化}的逻辑；注意两个绝对值的式子相加减，想一下三角不等式
        \end{jieda}
        \item 若实数$a$使得不等式$|x-2a| + |2x-a| \geq a^2$对任意的$x$恒成立，实数$a$的取值范围是\blkda{$\left[-\dfrac{3}{2}, \dfrac{3}{2}\right]$}
        \begin{jieda}
            \begin{enumerate}
            \item 方法1 \\
            首先将问题转化为左侧的最小值大于等于右侧，设函数$f(x) = |x-2a| + |2x-a|$，根据斜底锅函数的图象可知，其最低点必然位于系数绝对值最大的一项取零的位置，即$x = \dfrac{a}{2}$，最小值为$\left|\dfrac{3a}{2}\right|$，由此不等式可以转化为$\left|\dfrac{3a}{2}\right| \geq a^2$，即$\left |\dfrac{3}{2} \right| \geq |a|$，故$a \in \left[-\dfrac{3}{2}, \dfrac{3}{2}\right]$
            \item 方法2 \\
            通过设$x = ka(k \in \mathbf{R})$来简化计算。\\
            根据$x = ka(k \in \mathbf{R})$，不等式可以化为$|ka - 2a| + |2ka - a| \geq a^2$，\\
            消去$|a|$，不等式变为$|k - 2| + |2k - 1| \geq |a|$。\\
            通过对$k$进行分类讨论，可以得到$|k - 2| + |2k - 1|$的取值范围为$[\dfrac{3}{2}, +\infty)$，故$|a| \leq \dfrac{3}{2}$，即$a \in \left[-\dfrac{3}{2}, \dfrac{3}{2}\right]$
            \end{enumerate}
            \ps 斜底锅函数$f(x) = |ax+b|+|cx+d|$（$|a| \neq |c|$），若$|a| > |c|$，则取最小值时$x = -\dfrac{b}{a}$，若$|a| < |c|$，则取最小值时$x = -\dfrac{d}{c}$
        \end{jieda}
    \item 
    \begin{enumerate}
        \item  若关于$x$的不等式$\abs{x-1}-\abs{x-2}<a$的解集为$\bR$, 则实数$a$的取值范围为\blkda{$a>1$}
        
        \item 若关于$x$的不等式$\abs{x-1}-\abs{x-2}<a$有实数解, 则实数$a$的取值范围为\blkda{$a>-1$}. 
        
        \item 若关于$x$的不等式$\abs{x-1}-\abs{x-2}<a$的解集为$\varnothing$, 则实数$a$的取值范围为\blkda{$a\leq-1$}. 
    \end{enumerate}
    \begin{jieda}
        斜坡函数$f(x) = |x+a|-|x+b|$，最大值为$|a-b|$，最小值为$-|a-b|$. \\
        例如：$f(x) = |x+1|-|x+2| = \left \{ \begin{array}{ll}
            1 & x < -2 \\
            -2x-3 & x \in [-2, -1] \\
            -1 & x > -1 \\
        \end{array}\right.$
    \end{jieda}
    \end{enumerate}
    \item 基本不等式的应用
    \begin{enumerate}
        \item 
        \begin{enumerate}
            \item 若$a > 0, b > 0$，则$\dfrac{b}{a}+\dfrac{a}{b}$的取值范围是\blkda{$[2, +\infty)$}
            \item 若$ab > 0$，则$\dfrac{b}{a}+\dfrac{a}{b}$的取值范围是\blkda{$[2, +\infty)$}
            \item 若$ab < 0$，则$\dfrac{b}{a}+\dfrac{a}{b}$的取值范围是\blkda{$(-\infty, -2]$}
            \item 若$x > 0$，$x + \dfrac{3}{x}$的最小值是\blkda{$2\sqrt{3}$}
            \item 若$x > -1$，$x + \dfrac{3}{x+1}$的最小值是\blkda{$2\sqrt{3}-1$}
            \item 若$x \geq 0$，$x + \dfrac{3}{x+3}$的最小值是\blkda{$1$}
        \end{enumerate}
        \item 下列各式中, 最小值为4的是\dotfill(\kh{B})
        \xx
        {$x+2\sqrt{x}+5$}
        {$\dfrac{x^2+5}{\sqrt{x^2+1}}$}
        {$x^2+3+\dfrac{4}{x^2+3}$}
        {$x+\dfrac{4}{x}$}
        \begin{jieda}
            注意取等条件是否可以取到！
        \end{jieda}
        \item 已知不等式$(x+y)\left(\dfrac{a}{x}+\dfrac{1}{y}\right)\geq 9$对任意正实数$x,y$恒成立, 则正实数$a$的最小值是\dotfill(\kh{B})
        \xx
        {2}
        {4}
        {6}
        {8}
        \begin{jieda}
            首先求出$(x+y)\left(\dfrac{a}{x}+\dfrac{1}{y}\right)$最小值(含$a$)，然后令其大于等于9，得到$a$的取值范围，从中选取最小值\\
            $(x+y)\left(\dfrac{a}{x}+\dfrac{1}{y}\right) = a + 1 + \dfrac{ay}{x} + \dfrac{x}{y} \geq a + 1 + 2\sqrt{a}$（$x  = \sqrt{a}y$时取等）\\
            设$f(a) = a+1+2\sqrt{a}$，显然$f(a)$单调递增，且有$f(4) = 9$. 故$a \geq 4$时原式$\geq 9$，因此$a$最小值为4.
        \end{jieda}
    \item 已知关于$x$的不等式$(kx - k^2-4)(x-4) > 0$，其中$k \in \mathbf{R}$
    \begin{enumerate}
        \item 当$k$变化时，试求不等式的解集$A$
        \item 对于不等式的解集$A$，若满足$A \cap \mathbf{Z} = B$试探究集合$B$是否为有限集？若能，求出使集合$B$中元素个数最少的$k$的所有取值，并用列举法表示集合$B$；若不能，请说明理由。
    \end{enumerate}
    \begin{jieda}
        \begin{enumerate}
            \item 易知方程$(kx - k^2-4)(x-4) = 0$的两个解为$x_1 = 4, x_2 = k + \dfrac{4}{k}$（$k \neq 0$）
            \begin{dingautolist}{172}
                \item $k > 0$时，$k + \dfrac{4}{k} \geq 4$，原不等式的解集为$(-\infty, 4) \cup (k + \dfrac{4}{k}, +\infty)$
                \item $k = 0$时，原不等式变为$x-4 < 0$，因此解集为$(-\infty, 4)$
                \item $k < 0$时，$k + \dfrac{4}{k} < 4$，原不等式的解集为$(k + \dfrac{4}{k}, 4)$            \end{dingautolist}
            \item 根据前一问可知必然在$k < 0$时，$B = A \cap \mathbf{Z}$为有限集。当$k = -2$时，$k + \dfrac{4}{k}$取得最大值，此时$B = \{-3, -2, -1, 0, 1, 2, 3\}$
        \end{enumerate}
    \end{jieda}
    \item 若$a>b>0$, 则$a^2+\dfrac{16}{b(a-b)}$的最小值是\blkda{16}. 
    \begin{jieda}
        先考虑消元，$b(a-b) \leq \left (\dfrac{b+(a-b)}{2}\right )^2$，取等条件为$a = 2b$. \\
        故$a^2+\dfrac{16}{b(a-b)} \geq a^2 + \dfrac{64}{a^2} \geq 16$. 取等条件为$a = 2\sqrt{2}, b = \sqrt{2}$
    \end{jieda}
    \end{enumerate}
    \item 系数拆分求分式型代数式的最值
    \begin{enumerate}
        \item 已知$x, y, z \in \mathbf{R^+}$，则$\dfrac{xy+2yz}{x^2+y^2+z^2}$的最大值为\blkda{$\dfrac{\sqrt{5}}{2}$}
        \begin{jieda}
            \begin{enumerate}
                \item 首先观察可知$\dfrac{xy+2yz}{x^2+y^2+z^2}$的分子是乘积相加的形式，分母是平方相加的形式，因此应考虑采用待定系数法，利用平均值不等式将分子分母的形式统一。\\
                设$a>0, b > 0, a+b = 1, \dfrac{xy+2yz}{x^2+y^2+z^2} = \dfrac{xy+2yz}{x^2+ay^2+by^2+z^2} \leq \dfrac{xy+2yz}{2\sqrt{a}xy + 2\sqrt{b}yz}$ \\
                当$a = 0.2, b = 0.8$时，分子分母可以消去乘积项，变成常数。\\
                即$\dfrac{xy+2yz}{x^2+y^2+z^2} = \dfrac{xy+2yz}{x^2+0.2y^2+0.8y^2+z^2} \leq \dfrac{1}{2}\dfrac{xy+2yz}{\sqrt{0.2}xy + 2\sqrt{0.2}yz} = \dfrac{1}{2\sqrt{0.2}} = \dfrac{\sqrt{5}}{2}$
                \item $\dfrac{xy+2yz}{x^2+y^2+z^2} = \dfrac{\dfrac{x}{y}+\dfrac{2z}{y}}{1+\left(\dfrac{x}{y}\right)^2 + \left(\dfrac{z}{y}\right)^2}$

                设$\left\{\begin{array}{l}
                     \dfrac{x}{y} = rcos\theta  \\
                     \dfrac{z}{y} = rsin\theta
                \end{array}\right.$，则$\dfrac{\dfrac{x}{y}+\dfrac{2z}{y}}{1+\left(\dfrac{x}{y}\right)^2 + \left(\dfrac{z}{y}\right)^2} = \dfrac{r(cos\theta+2sin\theta)}{1+r^2} = \dfrac{cos\theta+2sin\theta}{r+\dfrac{1}{r}} \leq \dfrac{\sqrt{5}}{2}$
                \item 设$\left\{\begin{array}{l}
                     x = rcos\varphi cos\theta \\
                     y = rsin\varphi \\
                     z = rcos\varphi sin\theta 
                \end{array}\right.$，则$\dfrac{xy+2yz}{x^2+y^2+z^2} = \dfrac{r^2sin\varphi cos\varphi cos\theta + 2r^2sin\varphi cos\varphi sin\theta}{r^2} = sin\varphi cos\varphi cos\theta + 2sin\varphi cos\varphi sin\theta = \dfrac{sin2\varphi}{2} (cos\theta + 2sin\theta) \leq \dfrac{\sqrt{5}}{2}$
            \end{enumerate}
        \end{jieda}
        \item 若$x,y,z \in \mathbf{R^+}$，则$\dfrac{2xy + yz}{4x^2+4y^2+3z^2}$的最大值为\blkda{$\dfrac{\sqrt{3}}{6}$}
        \begin{jieda}
            分子的形式是乘积的和的形式，分母是平方和，易知分子分母之间可以用平均值不等式得到不等关系。观察可知分子中的$y$是复用的，因此考虑将分母中的$4y^2$拆成两部分分别与$x^2$和$z^2$配对。\\
            因此设$a+b=4$，$\dfrac{2xy+yz}{4x^2+ay^2+by^2+3z^2}$，根据平均值不等式有$\dfrac{2xy+yz}{4x^2+ay^2+by^2+3z^2} \leq \dfrac{2xy+yz}{4\sqrt{a}xy + 2\sqrt{3b}yz}$。\\
            当$4\sqrt{3b} = 4\sqrt{a}$的时候，不等号右侧为常数。此时$a = 3, b = 1$，对应的常数为$\dfrac{\sqrt{3}}{6}$，因为平均值不等式是恒成立的不等关系，因此原式最大值为$\dfrac{\sqrt{3}}{6}$
        \end{jieda}
    \end{enumerate}
    \item 有限制条件利用平均值不等式求最值-1的代换
    \begin{enumerate}
        \item 已知两个正实数$x, y$，$x+2y = 4$，求下列各式的最小值: \\
        \begin{enumerate*}
            \item $\dfrac{1}{x} + \dfrac{2}{y}$
            \item $\dfrac{1}{x+1} + \dfrac{2}{y+3}$
            \item $\dfrac{1}{2x+y} + \dfrac{2}{x+3y}$
        \end{enumerate*}
        \begin{jieda}
            \begin{enumerate}
                \item $\dfrac{1}{x} + \dfrac{2}{y} = (\dfrac{1}{x} + \dfrac{2}{y}) \cdot (x+2y) \cdot \dfrac{1}{4} = \dfrac{1}{4} \cdot (1 + 4 + \dfrac{2x}{y} + \dfrac{2y}{x}) \geq \dfrac{1}{4}(5+2\sqrt{4}) = \dfrac{9}{4}$（当$x = y = \dfrac{4}{3}$时取得最小值）
                \item $\dfrac{1}{x+1} + \dfrac{2}{y+3} = (\dfrac{1}{x+1} + \dfrac{2}{y+3}) \cdot (x+1+2y+6)\cdot\dfrac{1}{11} = \dfrac{1}{11}(1 + 4 + \dfrac{2x+2}{y+3} + \dfrac{2y+6}{x+1}) \\
                \geq \dfrac{1}{11}(5 + 2\sqrt{4}) = \dfrac{9}{11}$（当$x = \dfrac{8}{3}, y = \dfrac{2}{3}$时取得最小值）
                \item 设$m(2x+y) + n(x+3y) = k(x+2y)$，故$2(2m+n) = m+3n$，即$3m = n$，故可取$m = 1, n = 3$ \\
                $\dfrac{1}{2x+y} + \dfrac{2}{x+3y} = \dfrac{1}{2x+y} + \dfrac{6}{3x+9y} = (\dfrac{1}{2x+y} + \dfrac{6}{3x+9y}) \cdot (2x+y + 3x+9y) \cdot \dfrac{1}{20} \\
                = \dfrac{1}{20}(1 + 6 + \dfrac{3x+9y}{2x+y} + \dfrac{6(2x+y)}{3x+9y}) \geq \dfrac{1}{20}(7 + 2\sqrt{6}) = \dfrac{7+2\sqrt{6}}{20}$（当$\dfrac{3x+9y}{2x+y} = \dfrac{6(2x+y)}{3x+9y}$时取等）\\
                \ps 体会“基”的思想：将分母看做一组“基”，线性组合表示限制条件，根据系数对目标式进行变形.
            \end{enumerate}
        \end{jieda}
        \item 已知正实数$a, b$满足$a+b = 4$，则$\dfrac{1}{a+1} + \dfrac{1}{b+3}$的最小值是\blkda{$\dfrac{1}{2}$}.
        \begin{jieda}
            $\dfrac{1}{a+1} + \dfrac{1}{b+3} = (\dfrac{1}{a+1} + \dfrac{1}{b+3}) \times (a+1+b+3) \times \dfrac{1}{8} = (2+\dfrac{b+3}{a+1} + \dfrac{a+1}{b+3}) \times \dfrac{1}{8} \geq \dfrac{1}{2}$，当且仅当$a = 3, b = 1$时取等号.
        \end{jieda}
        \item 已知正实数$a, b$满足$4a+3b = 1$，则$\dfrac{1}{2a+b} + \dfrac{1}{a+b}$的最小值是\blkda{$3+2\sqrt{2}$}.
        \begin{jieda}
            $\dfrac{1}{2a+b} + \dfrac{1}{a+b} = \dfrac{1}{2a+b} + \dfrac{2}{2a+2b} = (\dfrac{1}{2a+b} + \dfrac{2}{2a+2b}) \times (2a + b + 2a + 2b) = 1 + 2 + \dfrac{2a+2b}{a+b} + \dfrac{2(2a+b)}{2a+2b} \geq 3+2\sqrt{2}$，当且仅当$a = \dfrac{3\sqrt{2}}{2}-2, b = 3-2\sqrt{2}$时取等号. 
        \end{jieda}
        \item 已知正实数$m, n$满足$\dfrac{1}{m+2n} + \dfrac{1}{2m+n} = 1$，则$m+n$的最小值为\blkda{$\dfrac{4}{3}$}
        \begin{jieda}
            $m+n = \dfrac{(m+2n) + (2m+n)}{3} = \dfrac{1}{3}[(m+2n) + (2m+n)] \cdot(\dfrac{1}{m+2n} + \dfrac{1}{2m+n}) \\
            = \dfrac{1}{3}(2 + \dfrac{2m+n}{m+2n} + \dfrac{m+2n}{2m+n}) \geq \dfrac{1}{3}(2 + 2\sqrt{1}) = \dfrac{4}{3}$（当$m = n = \dfrac{2}{3}$时取等号）
        \end{jieda}
    \end{enumerate}
    \item 有限制条件利用平均值不等式求最值-化为对勾
    \begin{enumerate}
        \item 已知$x>0, y>0, x+3y+xy=9$，则$x+3y$的最小值为\blkda{6}
        \begin{jieda}
            $x = \dfrac{9-3y}{1+y}$，故$x+3y = \dfrac{9-3y}{1+y} + 3y = 3y - 3 + \dfrac{12}{1+y} = 3(1+y) + \dfrac{12}{1+y} - 6 \\
            \geq 2\sqrt{3(1+y) \times \dfrac{12}{1+y}} - 6 = 6$
        \end{jieda}
        \item 已知$x > 0, y > 0, 4x+2y-xy=0$，则$2x+y$的最小值为\blkda{16}
        \begin{jieda}
            \begin{enumerate}
                \item 消元法 \\
                首先根据题目中的等式用$x$来表示$y$，得到$y = \dfrac{4x}{x-2} (x > 2)$ \\
                目标式子变为$2x+\dfrac{4x}{x-2} = 2x + 4 + \dfrac{8}{x-2} = 8 + 2(x-2) + \dfrac{8}{x-2} \geq 8 + 2 \sqrt{2(x-2) \cdot \dfrac{8}{x-2}} \\ = 16$
                \item 常数法 \\
                首先对题目中的等式进行变换，等式两边同时除以$xy$得到$\dfrac{2}{x} + \dfrac{4}{y} = 1$，所以有$(2x+y) = (2x+y)(\dfrac{2}{x} + \dfrac{4}{y}) = \dfrac{8x}{y} + \dfrac{2y}{x} + 8 \geq 2\sqrt{\dfrac{8x}{y} \cdot \dfrac{2y}{x}} + 8 = 16$
            \end{enumerate}
        \end{jieda}
        \item 已知$a > 0, b > 0$，且$ab = 1$，则$\dfrac{1}{2a} + \dfrac{1}{2b} + \dfrac{8}{a+b}$的最小值为\blkda{4}.
        \begin{jieda}
            \begin{enumerate}
                \item 消元法 \\
                用$a$表示$b$有$b = \dfrac{1}{a}$。原式可以化为$\dfrac{a}{2} + \dfrac{1}{2a} + \dfrac{8}{\dfrac{1}{a} + a} = \dfrac{1+a^2}{2a} + \dfrac{8a}{1+a^2} \geq 2\sqrt{\dfrac{1+a^2}{2a} \cdot \dfrac{8a}{1+a^2}} = 4$，当且仅当$a+b = a + \dfrac{1}{a} = 4$时取等号。 
                \item 常数法 \\
                原式$=(\dfrac{1}{2a} + \dfrac{1}{2b}) \times ab + \dfrac{8}{a+b} = \dfrac{a+b}{2} + \dfrac{8}{a+b} \geq 2\sqrt{\dfrac{a+b}{2} \cdot \dfrac{8}{a+b}} = 4$，当且仅当$\dfrac{a+b}{2} = \dfrac{8}{a+b}$，即$a+b=4$时取等号。
            \end{enumerate}
            \ps 容易出现的错解：消元后分别利用均值不等式得到答案5。\textbf{应用多次平均值不等式时应注意不等式方向是否相同，取等条件是否一致。}
        \end{jieda}
    \end{enumerate}
    \item 应用平均值不等式链求最值
    \begin{enumerate}
        \item $p: a>0, b>0$，$q_1:\sqrt{\dfrac{a^2+b^2}{2}} \geq \dfrac{a+b}{2}$，$q_2:\dfrac{a+b}{2} \geq \sqrt{ab}$，$q_3:\sqrt{ab} \geq \dfrac{2ab}{a+b}$，则\dotfill{（\kh{D}）}
        \xx{$p$是$q_1$的充要条件}{$p$是$q_2$的充要条件}{$p$是$q_3$的充要条件}{以上均不正确}
        \begin{jieda}
            $q_1$对任意$a, b$均成立，$q_2$对$a \geq 0, b \geq 0$均成立，$q_3$对$ab \geq 0$均成立，故$p$对$q_1, q_2, q_3$三者均是充分不必要条件.
        \end{jieda}
        \item 已知$a,b$为互不相等的正实数，则下列四个式子中最大的是 \dotfill{(\kh{B})}
        \xx{$\dfrac{2}{a+b}$}{$\dfrac{1}{a} + \dfrac{1}{b}$}{$\dfrac{2}{\sqrt{ab}}$}{$\sqrt{\dfrac{2}{a^2+b^2}}$}
        \begin{jieda}
            \begin{dingautolist}{172}
                \item 算术平均值大于几何平均值：$\dfrac{1}{a} + \dfrac{1}{b} > 2\sqrt{\dfrac{1}{a} \cdot \dfrac{1}{b}} = \dfrac{2}{\sqrt{ab}}$
                \item 算术平均值大于几何平均值：$\dfrac{2}{a+b} < \dfrac{2}{2\sqrt{ab}} < \dfrac{2}{\sqrt{ab}}$
                \item 算术平均值小于平方平均值：$\dfrac{2}{\sqrt{ab}} > \dfrac{2}{\sqrt{\frac{a^2+b^2}{2}}} = 2\sqrt{\dfrac{2}{a^2+b^2}} > \sqrt{\dfrac{2}{a^2+b^2}}$
            \end{dingautolist}
        \end{jieda}
        \item 设$a>0, b>0, a+b = 1$，则下列结论不正确的是\dotfill{(\kh{D})}
        \xx{$ab$的最大值为$\dfrac{1}{4}$}{$a^2+b^2$的最小值为$\dfrac{1}{2}$}{$\dfrac{4}{a} + \dfrac{1}{b}$的最小值为9}{$\sqrt{a}+\sqrt{b}$的最小值为$\sqrt{3}$}
        \begin{jieda}
            \begin{enumerate}[A.]
                \item $\sqrt{ab} \leq \dfrac{a+b}{2} \Longrightarrow ab \leq \dfrac{1}{4}$，（$a = b = \dfrac{1}{2}$时取等）
                \item $\sqrt{\dfrac{a^2+b^2}{2}} \geq \dfrac{a+b}{2} \Longrightarrow a^2+b^2 \geq \dfrac{1}{2}$（$a = b = \dfrac{1}{2}$时取等）
                \item $\dfrac{4}{a} + \dfrac{1}{b} = (\dfrac{4}{a} + \dfrac{1}{b}) \cdot (a+b) = 5 + \dfrac{4b}{a} + \dfrac{a}{b} \geq 5 + 2\sqrt{4} = 9$（$a = \dfrac{2}{3}, b = \dfrac{1}{3}$时取等）
                \item $\dfrac{\sqrt{a} + \sqrt{b}}{2} \leq \sqrt{\dfrac{a+b}{2}}$，故$\sqrt{a} + \sqrt{b} \leq \sqrt{2}$
            \end{enumerate}
        \end{jieda}
    \end{enumerate}
    \item 不等式恒成立问题
    \begin{enumerate}
        \item 已知不等式$x^2+(a-2)x+3a>0(*)$
        \begin{enumerate}
            \item 若不等式$(*)$对任意的$x \in [-1, 1]$恒成立，求实数$a$的取值范围
            \item 若不等式$(*)$对任意的$a \in [-1, 1]$恒成立，求实数$x$的取值范围
        \end{enumerate}
        \begin{jieda}
            \begin{enumerate}
                \item 
                \begin{dingautolist}{172}
                    \item 通解法 \\
                    $f(x) = x^2+(a-2)x+3a$的对称轴为$x = 1-\dfrac{a}{2}$ \\
                    若$1-\dfrac{a}{2} < -1$，即$a > 4$，此时$f(x)$在$[-1, 1]$上单调递增，因此$f(-1) > 0$，即$1+2-a+3a > 0$，解得$a > -\dfrac{3}{2}$，结合对称轴有$a > 4$ \\
                    若$1-\dfrac{a}{2} > 1$，即$a < 0$，此时$f(x)$在$[-1, 1]$上单调递减，因此$f(1) > 0$，即$1+a-2+3a > 0$，解得$a > \dfrac{1}{4}$，此种情况无解 \\
                    若$1-\dfrac{a}{2} \in [-1, 1]$，即$a \in [0, 4]$，此时$f(x)$在$[-1, 1]$内最小值在顶点处取，即$f(1-\dfrac{a}{2}) > 0$，整理有$a^2-16a+4<0$，解得$a \in (8-2\sqrt{15}, 8+2\sqrt{15})$，结合对称轴有$a \in (8-2\sqrt{15}, 4]$ \\
                    综上，$a \in (8-2\sqrt{15}, +\infty)$
                    \item 参变分离 \\
                    $(x+3)a-2x+x^2 > 0$在$x \in [-1, 1]$恒成立，即$a > \dfrac{2x-x^2}{x+3}$在$x \in [-1, 1]$恒成立.\\
                    换元$t = x+3$， $t \in [2, 4]$，即$a > \dfrac{-t^2+8t-15}{t}$在$[2, 4]$恒成立.进一步变形有$a > -(t+\dfrac{15}{t})+8$在$[2, 4]$恒成立.根据对勾函数性质可知$-(t+\dfrac{15}{t})+8 \leq 8 - 2\sqrt{15}$（$t = \sqrt{15}$，即$x = \sqrt{15}-3$时取等号）,故$a \in (8-2\sqrt{15}, +\infty)$
                \end{dingautolist}
                \item 将$a$做为主元，即$(x+3)a-2x+x^2 > 0$在$a \in [-1, 1]$恒成立，设$f(a) = (x+3)a-2x+x^2$，故只需$f(1) > 0, f(-1) > 0$即可.解得$x \in \left(-\infty, \dfrac{3-\sqrt{21}}{2}\right)\cup\left(\dfrac{3+\sqrt{21}}{2}, +\infty\right)$
            \end{enumerate}
            \ps 恒成立问题，首先考虑主元应该是谁，然后考虑是否要参变分离.
        \end{jieda}
        \item 若不等式$x^2+px>4x+p-3$，当$0 \leq p \leq 4$时恒成立，则$x$的取值范围是\blkda[4]{$\left(-\infty, -1\right) \cup \left(3, +\infty\right)$}
        \begin{jieda}
            \begin{enumerate}
                \item 方法1：看做是关于$p$的一元一次不等式求解 \\
                不等式移项得到$x^2 -4x + 3 > (1-x)p$。\\
                因式分解有$(x-1)(x-3) > -p(x-1)$ \\
                下面根据$x$的取值范围进行分类讨论：
                \begin{dingautolist}{172}
                    \item $x > 1$，不等式转化为$x - 3 > -p$，故$x \in (3, +\infty)$
                    \item $x = 1$，这种情况不成立
                    \item $x < 1$，不等式转化为$x - 3 < -p$，故$x \in (-\infty, -1)$
                \end{dingautolist}
                综上，当$x \in (-\infty, -1) \cup (3, +\infty)$时不等式恒成立。
                \item 方法2：看做是关于$x$的一元二次不等式求解 \\
                不等式移项得到$x^2 + (p-4)x + (3-p) > 0$。即$(x-(3-p))(x-1) > 0$ \\
                因此根据两解可以写出不等式的解集为$\left \{
                    \begin{array}{ll}
                         (-\infty, 3-p) \cup (1, +\infty) & 4 \geq p \geq 2 \\
                         (-\infty, 1) \cup (3-p, +\infty) & 0 \leq p < 2
                    \end{array}
                \right. $ \\
                要使得不等式在$p$取到$[0, 4]$内任意值时成立，则应取所有$p$所对应的解集的交集，故当$x \in (-\infty, -1) \cup (3, +\infty)$时不等式恒成立。
            \end{enumerate}
            \ps 通过对比可知以$p$为主元更容易理解和推广。
        \end{jieda}
        \item 关于$x$的不等式$x^2+9+|x^2-3x| \geq kx$在区间$[1,5]$上恒成立，则实数$k$的取值范围是\blkda{$k \leq 6$}
        \begin{jieda}
            当$x \in [1, 5]$时，原不等式可以转化为$x + \dfrac{9}{x} + |x - 3| \geq k$ \\
            不等式恒成立即$k$比左侧最小值小，而易知左侧最小值在$x = 3$时取，故$k \leq 6$
        \end{jieda}
        \item 若不等式$\dfrac{2x^2+2kx+k}{4x^2+6x+3} < 1$对一切$x \in \mathbf{R}$恒成立，则$k$的取值范围是\blkda{$(1,3)$}
        \begin{jieda}
            首先发现$4x^2+6x+3 > 0$恒成立，因此不等式转化为$2x^2+2kx+k < 4x^2+6x+3$，即$2x^2 + (6-2k)x + (3-k) > 0$恒成立，故有$\Delta = (6-2k)^2 - 8(3-k) = 4k^2 -16k + 12 < 0$，解得$k \in (1,3)$.
        \end{jieda}
        \item 对任意的$x \in \mathbf{R}$，不等式$\dfrac{3x^2 + 2x + 2}{x^2+x+1} > k$恒成立，求正整数$k$的取值范围.
        \begin{jieda}
            首先发现$x^2+x+1 > 0$恒成立，因此原不等式转化为$(3-k)x^2 + (2-k)x + (2-k) > 0$。此处应当分类讨论
            \begin{enumerate}
                \item $k = 3$，此时退化为一次形式，不等关系不恒成立，不符合题意
                \item $k \neq 3$，此时对应的二次函数应当开口向上，所以有$k < 3$，计算判别式小于0可得$k < 2$，因此$k = 1$
            \end{enumerate}
            综上，$k \in \{1\}$
        \end{jieda}
        \item *已知实数$m$满足：当关于$x$的实系数一元二次方程$ax^2+bx+c=0$有实数根时，$(a-b)^2+(b-c)^2+(c-a)^2 \geq ma^2$恒成立，求$m$的最大值.
        \begin{jieda}
            参变分离有：$\dfrac{(a-b)^2}{a^2}+\dfrac{(b-c)^2}{a^2}+\dfrac{(c-a)^2}{a^2} \geq m$，即$(1-\dfrac{b}{a})^2 + (\dfrac{b}{a} - \dfrac{c}{a})^2 + (\dfrac{c}{a} - 1)^2 \geq m$。设方程$ax^2+bx+c=0$的两根为$x_1, x_2$，故根据韦达定理有$(1+x_1+x_2)^2 + (x_1+x_2+x_1x_2)^2 + (x_1x_2-1)^2 \geq m$。\\
            展开得$1+x_1^2+x_2^2+2x_1 + 2x_2+2x_1x_2 + x_1^2 + x_2^2 + x_1^2x_2^2 + 2x_1x_2 + 2x_1^2x_2 + 2x_1x_2^2 + x_1^2x_2^2 -2x_1x_2 + 1 \geq m$
            整理得$2(x_1^2x_2^2 + x_1^2x_2 + x_1x_2^2 + x_1^2 + x_2^2 + x_1x_2 + x_1 + x_2 + 1) \geq m$ \\
            观察式子可知$x_1, x_2$对称，共9项，故可尝试因式分解。得到$2(x_1^2 + x_1 + 1)(x_2^2 + x_2 + 1) \geq m$。可以看成以$x_1, x_2$为变量的两个二次函数相乘，该二次函数当自变量取$-\dfrac{1}{2}$时取最小值$\dfrac{3}{4}$，因此$m \leq 2 \times \dfrac{3}{4} \times \dfrac{3}{4} = \dfrac{9}{8}$
        \end{jieda} 
    \end{enumerate}
    \item 不等式有解问题
    \begin{enumerate}
        \item 关于$x$的不等式$x^2+9+|x^2-3x| \geq kx$在区间$[1,5]$上有解，则实数$k$的取值范围是\blkda{$k \leq 12$}
        \begin{jieda}
            当$x \in [1, 5]$时，原不等式可以转化为$x + \dfrac{9}{x} + |x - 3| \geq k$ \\
            不等式有解即$k$比左侧最大值小，设$f(x) = x + \dfrac{9}{x} + |x - 3|$易知$f(x)$在$[1, 5]$上先减后增，最大值在端点处取. $f(1) = 12, f(5) = \dfrac{44}{5}$故$k \leq 12$\\
            \ps \textbf{区分有解问题和恒成立问题}.
        \end{jieda}
        \item 已知函数$f(x) = \left|x + \dfrac{1}{x} - 4\right|$，若关于$x$的不等式$f(x) \geq m^2-m+2$在区间$\left[\dfrac{1}{6}, 3\right]$总有解，则实数$m$的取值范围是\blkda[3.5]{$\left[\dfrac{3-\sqrt{15}}{6}, \dfrac{3+\sqrt{15}}{6}\right]$}
        \begin{jieda}
            首先要意识到不等式的右侧没有$x$，因此可以将右侧当做一个整体，直接找$f(x)$在区间$\left[\dfrac{1}{6}, 3\right]$上的最大值. 根据对勾函数的性质可知$f(x)$在$\left[\dfrac{1}{6}, 1\right]$上先减后增，在$[1, 3]$单调递减. 故最大值会在$x = \dfrac{1}{6}$或$x = 1$处取到.$f\left(\dfrac{1}{6}\right) = \dfrac{13}{6}, f(1) = 2$，故$f(x)$的最大值为$\dfrac{13}{6}$ \\
            问题转化为求解不等式$\dfrac{13}{6} \geq m^2-m+2$，解集为$\left[\dfrac{3-\sqrt{15}}{6}, \dfrac{3+\sqrt{15}}{6}\right]$. 
        \end{jieda}    
    \end{enumerate}
    \item 命题结合一元二次方程、一元二次不等式的有解/恒成立问题
    \begin{enumerate}
        \item 若“$\exists x \in \mathbf{R}, x^2 - 4x +a = 0$”为真命题，求实数$a$的取值范围．
        \begin{jieda}
            首先将题干信息进行转化，“$\exists x \in \mathbf{R}, x^2 - 4x + a = 0$”为真命题即方程$x^2 - 4x + a = 0$有实根，因此其判别式$\Delta = 16 - 4a \geq 0$，故$a \leq 4$
        \end{jieda}
        \item 若“$\forall x \in \mathbf{R}, x^2 + 3ax+ 9 > 0$”，“$\exists x \in \mathbf{R}, x^2 + 2x + a < 0$”均为真命题，则$a$的取值范围是\blkda{$(-2, 1)$}
        \begin{jieda}
            “$\forall x \in \mathbf{R}, x^2 + 3ax+ 9 > 0$”是真命题，即方程$x^2 + 3ax + 9 = 0$的判别式$\Delta = 9a^2 - 36 < 0$，解得$a \in (-2, 2)$\\
            “$\exists x \in \mathbf{R}, x^2 + 2x + a < 0$”是真命题，即方程$x^2 + 2x + a = 0$的判别式$\Delta = 4 - 4a > 0$，解得$a \in (-\infty, 1)$ \\
            两者求交集，得到$a \in (-2, 1)$
        \end{jieda}
        \item 若“$\exists x \in \mathbf{R}, x^2 - 2mx + m + 2 < 0$”为真命题，则实数$m$的取值范围是\blkda[4]{$(-\infty， -1) \cup (2, +\infty)$}
        \begin{jieda}
            “$\exists x \in \mathbf{R}, x^2 - 2mx + m + 2 < 0$”为真命题即方程$x^2 - 2mx + m + 2 = 0$的判别式$\Delta = 4m^2 - 4(m+2) > 0$，整理得$m^2 - m - 2 > 0$，解得$m \in (-\infty， -1) \cup (2, +\infty)$
        \end{jieda}
        \item 若“$\exists x \in [-2, 0]$，$m \geq x^2 +2x -3$”是真命题，则实数$m$的取值范围是\blkda{$[-4, +\infty)$}
        \begin{jieda}
            题目可以转化为$m \geq x^2 +2x -3$在$[-2, 0]$内的最小值。由二次函数开口向上，对称轴为$x = -1$可知最小值为$-4$，故$m \in [-4, +\infty)$
        \end{jieda}
        \item 若“$\forall x \in [1, 4], x^2 -ax + 1 \leq 0$”为真命题，则实数$a$的取值范围为\blkda{$\left[\dfrac{17}{4}, +\infty \right)$}
        \begin{jieda}
            \begin{enumerate}
                \item 设$f(x) = x^2 -ax + 1$，则“$\forall x \in [1, 4], x^2 -ax + 1 \leq 0$”为真命题可以转化为$f(x)$在$[1, 4]$的最大值小于等于0。易知$f(x)$开口向上，最大值在端点处取到，因此应有$2 - a \leq 0$且$17 - 4a \leq 0$，故$a \geq \dfrac{17}{4}$
                \item “$\forall x \in [1, 4], x^2 -ax + 1 \leq 0$”为真命题可以转化为$a \geq \dfrac{x^2+1}{x}$对于$x \in [1, 4]$恒成立，即$a \geq x + \dfrac{1}{x}$在$[1, 4]$上的最大值。根据对勾函数的性质可知其在$[1, 4]$上的最大值为$\dfrac{17}{4}$，故$a \geq \dfrac{17}{4}$
            \end{enumerate}
        \end{jieda}        
        \item 若“$\forall m \in [-1, 1], x^2 + (m-4)x + 4 - 2m > 0$”为真命题，则$x$的取值范围是\blkda[3]{$(-\infty, 1) \cup (3, +\infty)$}
        \begin{jieda}
            \begin{enumerate}
            \item 方法1 \\
            将不等式$x^2 + (m-4)x + 4 - 2m > 0$看做是关于关于$m$的一次不等式，因此考虑设$f(m) = (x-2)m + (x-2)^2$，故不等式$f(m) > 0$在$[-1, 1]$内恒成立可以转化为$f(1) > 0$且$f(-1) > 0$，解得$x \in (-\infty, 1) \cup (3, +\infty)$
            \item 方法2 \\
            首先可以利用因式分解或求根公式得到$x^2 + (m-4)x + 4 - 2m = 0$的两根为$2$和$2-m$。
            \begin{dingautolist}{172}
                \item 当$2 > 2-m$，即$m \in (0, 1]$时，$x$的取值范围是$(-\infty, 2-m) \cup (2, +\infty)$，对$\forall m \in (0, 1]$都成立，即对所有的结果取交集，即$x \in (-\infty, -1) \cup (2, +\infty)$
                \item 当$2-m > 2$时，即$m \in [-1, 0)$时，$x$的取值范围是$(-\infty, 2) \cup (2-m, +\infty)$，对$\forall m \in [-1, 0)$都成立，即对所有的结果取交集，即$x \in (-\infty, 2) \cup (3, +\infty)$
                \item 当$m = 0$时，$x$的取值范围是$(-\infty, 2) \cup (2, +\infty)$
            \end{dingautolist}
            因为要对所有的$m \in [-1, 1]$都成立，因此应对三种情况取交集，故有$x \in (-\infty, 1) \cup (3, +\infty)$
        \end{enumerate}
        \end{jieda}
    \end{enumerate}
    
    \item 二次方程根的分布
    \begin{enumerate}
        \item 已知$a, b$为常数，函数$f(x) = x^2 - bx + a$，当$a = 2b-1$时，若方程$f(x) = 0$在$(-2, 1)$上有解，求实数$b$的取值范围.
        \begin{jieda}
            $f(x) = x^2 - bx + 2b - 1$.
            \begin{dingautolist}{172}
                \item 分对称轴位置讨论 \\
                $f(x)$的对称轴为$x = \dfrac{b}{2}$，$f(-2) = 4b+3$，$f(1) = b$，$\Delta = b^2 - 4(2b-1) = b^2 - 8b + 4$
                \begin{enumerate}[1.]
                    \item 若$\dfrac{b}{2} \leq -2$，则$f(x)$在$(-2, 1)$单调增，故“方程$f(x) = 0$在$(-2, 1)$上有解”即$f(-2) < 0, f(1) > 0$，此种情况解集为空集.
                    \item 若$\dfrac{b}{2} \geq 1$，则$f(x)$在$(-2, 1)$单调减，故“方程$f(x) = 0$在$(-2, 1)$上有解”即$f(-2) > 0, f(1) < 0$，此种情况解集为空集.
                    \item 若$\dfrac{b}{2} \in (-2, 1)$，则$f(x)$在$(-2, 1)$先减后增，故“方程$f(x) = 0$在$(-2, 1)$上有解”即$\Delta \geq 0$且有$f(-2) > 0$或$f(1) > 0$，此种情况解集为$(-\dfrac{3}{4}, 4-2\sqrt{3}]$
                \end{enumerate}
                综上，$b \in (-\dfrac{3}{4}, 4-2\sqrt{3}]$
                \item 参变分离 \\
                “方程$f(x) = 0$在$(-2, 1)$上有解”即“$b = \dfrac{1-x^2}{2-x}$在$(-2, 1)$上有解”，即“$b = \dfrac{3}{x-2} + (x-2) + 4$在$(-2, 1)$上有解” \\
                因此$b$的取值范围即$y = \dfrac{3}{x-2} + (x-2) + 4$在$(-2, 1)$上的值域. \\
                $y = \dfrac{3}{x-2} + (x-2) + 4$在$(-2, \sqrt{3}-2)$单调递增，在$[2-\sqrt{3}, 1)$单调递减，值域为$(-\dfrac{3}{4}, 4-2\sqrt{3}]$，故$b \in (-\dfrac{3}{4}, 4-2\sqrt{3}]$
            \end{dingautolist}
            \ps 有解问题可以采用分对称轴位置讨论或参变分离解决，\textbf{小题推荐参变分离，大题可以分类讨论写步骤，参变分离检验}。分对称轴位置讨论时，对于对称轴在目标区间内的情况要注意其充要条件。参变分离会把方程有解问题转化成函数求值域问题。
        \end{jieda}
        \item 已知函数$f(x) = x^2 + 2mx + 2m + 3$至少有一个零点在区间$(0, 2)$内，求实数$m$的取值范围.    \begin{jieda}
        \begin{dingautolist}{172}
            \item 分对称轴位置讨论 \\
            $f(x)$对称轴为$x = -m$，$f(0) = 2m+3$，$f(2) = 7+6m$，$\Delta = 4m^2 - 8m - 12$
            \begin{enumerate}[1.]
                \item 若$-m \leq 0$，则$f(x)$在$(0, 2)$单调增，故“$f(x)$在$(0, 2)$上有零点”即$f(0) < 0, f(2) > 0$，此种情况无解.
                \item $-m \geq 2$，则$f(x)$在$(0, 2)$单调减，故“$f(x)$在$(0, 2)$上有零点”即$f(0) > 0, f(2) < 0$，此种情况无解.
                \item $-m \in (0, 2)$，则$f(x)$在$(0, 2)$先减后增，“$f(x)$在$(0, 2)$上有零点”即$\Delta \geq 0$，且有$f(0) > 0$或$f(2) > 0$，此种情况解得$m \in \left(-\dfrac{3}{2}, -1\right]$
            \end{enumerate}
            综上，$m \in \left(-\dfrac{3}{2}, -1\right]$
            \item 参变分离 \\
            $f(x)$至少有一个零点在区间$(0, 2)$内，即$f(x) = 0$在$(0, 2)$内有根，即方程$m = \dfrac{-3-x^2}{2x+2}$在$(0, 2)$内有根，即方程$m = \dfrac{1}{2} \cdot \dfrac{-3-(t-1)^2}{t}$在$t \in (1, 3)$内有根，即方程$m = \dfrac{1}{2} \cdot \left[-(t+\dfrac{4}{t})+2\right]$在$t \in (1, 3)$内有根.\\
            因此$m$的取值范围即$y = \dfrac{1}{2} \cdot [-(t+\dfrac{4}{t})+2]$在$t \in (1, 3)$内的值域. \\
            $y = \dfrac{1}{2} \cdot \left[-(t+\dfrac{4}{t})+2\right]$在$(1, 2)$单调增，在$(2, 3)$单调减，因此值域为$\left(-\dfrac{3}{2}, -1\right]$，故$m \in \left(-\dfrac{3}{2}, -1\right]$
        \end{dingautolist}
    \end{jieda}
        \item 已知函数$f(x) = x^2 + 2mx + 2m + 3$恰有一个零点在区间$(0, 2)$内，求实数$m$的取值范围.
        \begin{jieda}
            \begin{dingautolist}{172}
                \item 通解法 \\
                $f(x)$对称轴为$x = -m$，$f(0) = 2m+3$，$f(2) = 7+6m$，$\Delta = 4m^2 - 8m - 12$
                \begin{enumerate}[1.]
                    \item 符合函数零点存在定理，此时有$f(0) \cdot f(2) < 0$，即$m \in (-\dfrac{3}{2}, -\dfrac{7}{6})$
                    \item 二次函数与$x$轴相切在区间内，此时有$\Delta = 0$，且$-m \in (0, 2)$即$m = -1$
                    \item 左端点函数值为0，此时有$m = -\dfrac{3}{2}$，此时$f(2) < 0$，$f(x)$在$(0, 2)$内先减后增，故不存在零点.
                    \item 右端点函数值为0，此时有$m = -\dfrac{7}{6}$，此时$f(0) > 0$，$f(x)$在$(0, 2)$内先减后增，存在唯一零点.
                \end{enumerate}
                综上，$m \in (-\dfrac{3}{2}, -\dfrac{7}{6}] \cup \{-1\}$
                \item 参变分离 \\
                $f(x)$有一个零点在区间$(0, 2)$内，即方程$m = \dfrac{-3-x^2}{2x+2}$在$(0, 2)$内有一根，即方程$m = \dfrac{1}{2} \cdot \dfrac{-3-(t-1)^2}{t}$在$t \in (1, 3)$内有一根，即方程$m = \dfrac{1}{2} \cdot [-(t+\dfrac{4}{t})+2]$在$t \in (1, 3)$内有一根，即函数$y = m$与函数$y = \dfrac{1}{2} \cdot [-(t+\dfrac{4}{t})+2]$在$t \in (1, 3)$内只有一个交点.\\
                $y = \dfrac{1}{2} \cdot [-(t+\dfrac{4}{t})+2]$在$(1, 2)$单调增，在$(2, 3)$单调减，因此当$m \in (-\dfrac{3}{2}, -\dfrac{7}{6}] \cup \{-1\}$时满足要求.
            \end{dingautolist}
            \ps 二次函数在区间内只有一个零点的通解法对应着\textbf{四种情况}：\begin{enumerate*}[1.]
                \item 符合函数零点存在定理
                \item $\Delta = 0$且对称轴在区间内
                \item 左端点为函数的一个零点（应根据区间开闭情况单独检验）
                \item 右端点为函数的一个零点（应根据区间开闭情况单独检验）
            \end{enumerate*} \\
            参变分离会将问题转化为动直线和定曲线只有唯一交点，要\textbf{注意相切的情况和端点情况}。\\
            这种情况参变分离和通解法都没有明显优势.
        \end{jieda}
        \item 已知函数$f(x) = x^2 + 2mx + 2m + 3$恰有一个零点在区间$(0, +\infty)$内，求实数$m$的取值范围.
        \begin{jieda}
            \begin{dingautolist}{172}
                \item 通解法 \\
                $f(x)$对称轴为$x = -m$，$f(0) = 2m+3$，$\Delta = 4m^2 - 8m - 12$
                \begin{enumerate}[1.]
                    \item “符合函数零点存在定理”，此时有$f(0) < 0$，即$m \in (-\infty, -\dfrac{3}{2})$
                    \item 二次函数与$x$轴相切在区间内，此时有$\Delta = 0$，且$-m \in (0, +\infty)$即$m = -1$
                    \item 左端点函数值为0，此时有$m = -\dfrac{3}{2}$，$f(x)$在$(0, +\infty)$内先减后增，存在唯一零点.
                \end{enumerate}
                综上，$m \in (-\infty, -\dfrac{3}{2}] \cup \{-1\}$
                \item 参变分离 \\
                $f(x)$有一个零点在区间$(0, +\infty)$内，即方程$m = \dfrac{-3-x^2}{2x+2}$在$(0, +\infty)$内有一根，即方程$m = \dfrac{1}{2} \cdot \dfrac{-3-(t-1)^2}{t}$在$t \in (1, +\infty)$内有一根，即方程$m = \dfrac{1}{2} \cdot [-(t+\dfrac{4}{t})+2]$在$t \in (1, +\infty)$内有一根，即函数$y = m$与函数$y = \dfrac{1}{2} \cdot [-(t+\dfrac{4}{t})+2]$在$t \in (1, +\infty)$内只有一个交点.\\
                $y = \dfrac{1}{2} \cdot [-(t+\dfrac{4}{t})+2]$在$(1, 2)$单调增，在$(2, +\infty)$单调减，因此当$m \in (-\infty, -\dfrac{3}{2}] \cup \{-1\}$时满足要求.
            \end{dingautolist}
            \ps 区间一侧为无穷的情况对于通解法可以根据开口方向将问题简化.
        \end{jieda}
        \item 已知$g(x) = -x^2+2x$，若关于$x$的方程$g(x) = -mx-m$在$(0, 2)$上恰有两个不等实根，求$m$的取值范围.
        \begin{jieda}
            \begin{dingautolist}{172}
                \item 通解法 \\
                “方程$g(x) = -mx-m$在$(0, 2)$上恰有两个不等实根”即“方程$-x^2+(2+m)x+m = 0$在$(0, 2)$上恰有两个不等实根”，设$h(x) = -x^2+(2+m)x+m$，对称轴为$x = \dfrac{2+m}{2}$，$\Delta = (2+m)^2 + 4m$，$h(0) = m$，$h(2) = 3m$ \\
                $h(x) = 0$在$(0, 2)$上恰有两个不等实根的充要条件是$\left \{\begin{array}{l}
                    \Delta > 0  \\
                    \dfrac{2+m}{2} \in (0, 2) \\
                    h(0) < 0 \\
                    h(2) < 0
                \end{array} \right.$，解得$m \in (2\sqrt{3}-4, 0)$
                \item 参变分离 \\
                “方程$g(x) = -mx-m$在$(0, 2)$上恰有两个不等实根”即“方程$m = \dfrac{x^2-2x}{x+1}$在$(0, 2)$上恰有两个不等实根”，即“函数$y = m$与函数$y = t + \dfrac{3}{t} - 4$在$(1, 3)$上恰有两个交点” \\
                $y = t + \dfrac{3}{t} - 4$在$(1, \sqrt{3})$单调减，在$[\sqrt{3}, 3)$上单调增，因此当$m \in (2\sqrt{3}-4, 0)$时满足要求.
            \end{dingautolist}
        \end{jieda}
        \item 已知方程$x^2+2mx+m+2 = 0$有两个不相等的正实根，则实数$m$的取值范围是\blkda{$(-2, -1)$}
        \begin{jieda}
            设$f(x) = x^2+2mx+m+2$，则$f(x) = 0$有两个不相等的正实根即$\left \{ \begin{array}{l}
                 \Delta > 0 \\ 
                 f(0) > 0  \\
                 -m > 0
            \end{array}\right .$，即$\left \{ \begin{array}{l}
                 m \in (-\infty, -1) \cup (2, +\infty) \\ 
                 m > -2  \\
                 m < 0
            \end{array}\right .$，因此$m \in (-2, -1)$
        \end{jieda}
    \end{enumerate}
    \item 平均值不等式相关的综合应用
    \begin{enumerate}
        \item 已知两两不同的$x_1, y_1, x_2, y_2, x_3, y_3$满足$x_1 + y_1 = x_2 + y_2 = x_3 + y_3$，且$x_1 < y_1, \quad x_2 < y_2, \quad x_3 < y_3, \quad x_1y_1+x_3y_3=2x_2y_2 > 0$，则下列选项中恒成立的是 \dotfill(\kh{A})
        \xx{$2x_2 < x_1 + x_3$}{$2x_2 > x_1 + x_3$}{$x_2^2 < x_1x_3$}{$x_2^2 > x_1x_3$}
        \begin{jieda}
            设$x_1+y_1 = 2m$，故可设$\left \{
        \begin{array}{lll}
             x_1 = m - a, & y_1 = m + a, & a > 0\\
             x_2 = m - b, & y_2 = m + b, & b > 0\\
             x_3 = m - c, & y_3 = m + c, & c > 0
        \end{array}
        \right.$  \\
        由此，题目中的条件可以转化为 $m^2 - a^2 + m^2 - c^2 = 2(m^2-b^2)$，即$a^2+c^2 = 2b^2$，且有$m^2 - b^2 > 0$。下面推导$a+c<2b$
        \begin{enumerate}
            \item 方法1：利用平方差公式 \\
            原式移项后可以进一步可以转化为$(a+b)(a-b) + (c+b)(c-b) = 0$。\\
            $(a+b)(a-b) + (c+b)(c-b) = 0$说明$a < b < c$或者$c < b < a$，不妨设$a < b < c$（另一种情况对称），此时必有$(a+b) < (c+b)$，因此有$|a-b| > |c-b|$，即$b-a > c-b$，故$a+c<2b$
            \item 方法2：利用平均值不等式 \\
            由平均值不等式有$\sqrt{\dfrac{a^2+c^2}{2}} > \dfrac{a+c}{2}$，根据$a^2+c^2 = 2b^2$，可以转化为$2b > a + c$
        \end{enumerate}
        \end{jieda} 
        \item 已知$a,b,c,d$满足$a^2-ab+4 = 0, c^2+d^2 = 1$，则当$(a-c)^2 + (b-d)^2$取最小值的时候，$abcd = $\blkda{$1+\sqrt{2}$}
        \begin{jieda}
            $(a-c)^2 + (b-d)^2$的几何意义为$A(a, b), B(c,d)$两点间距离的平方，$A$在对勾函数$y = x + \dfrac{4}{x}$上，$B$在单位圆$x^2+y^2 = 1$上。$(a-c)^2 + (b-d)^2$为$|AB|^2$ \\
            固定点$A$时，$OAB$构成三角形$|AB| + 1 > |OB|$,仅当$B$在$OA$上时此时$|AB| = |OB| - 1$最小。 \\
            $|OA|^2 = a^2 + b^2 = a^2 + (a + \dfrac{4}{a})^2 = 2a^2 + 8 + \dfrac{16}{a^2} \geq 8 + 2\sqrt{32} = 8 + 8\sqrt{2}$（$a = \sqrt[4]{8}$时取等）\\
            此时$ab = a^2+4 = 4 + 2\sqrt{2}$，此时根据相似三角形可知$\dfrac{a}{c} = \dfrac{b}{d} = \dfrac{|OA|}{1} = \sqrt{8 + 8\sqrt{2}}$，故$\dfrac{ab}{cd} = 8 + 8\sqrt{2}$，因此$cd = \dfrac{4 + 2\sqrt{2}}{8+8\sqrt{2}}$，故$abcd = \dfrac{(4 + 2\sqrt{2})^2}{8+8\sqrt{2}} = \dfrac{24 + 16\sqrt{2}}{8+8\sqrt{2}} = \dfrac{3+2\sqrt{2}}{1+\sqrt{2}} = 1+\sqrt{2}$
        \end{jieda}
        \item 已知一元二次函数$f(x) = ax^2 + bx + c (a > 0, c > 0)$的图像与$x$轴有两个不同的公共点，其中一个公共点的坐标为$(c, 0)$，且当$0 < x < c$时，恒有$f(x) > 0$。
        \begin{enumerate}
            \item 当$a = 1, c = \dfrac{1}{2}$时，求出不等式$f(x) < 0$的解
            \item 求出不等式$f(x) < 0$的解（用$a, c$表示）
            \item 若以二次函数的图象与坐标轴的三个交点为顶点的三角形面积为8，求$a$的取值范围
            \item 若不等式$m^2-2km+1+b+ac \geq 0$对所有$k \in [-1, 1]$恒成立，求$m$的取值范围。
        \end{enumerate}
        \begin{jieda}
            \begin{enumerate}
                \item 当$a = 1, c = \dfrac{1}{2}$时，$f(x) = x^2 + bx + \dfrac{1}{2}$，因为过$(\dfrac{1}{2}, 0)$，所以$b = -\dfrac{3}{2}$，因此有$f(1) = 0$，因此$f(x) < 0$的解集为$(\dfrac{1}{2}, 1)$
                \item 根据韦达定理可知$f(\dfrac{1}{a}) = f(c) = 0$，因为当$0 < x < c$时，恒有$f(x) > 0$，所以必有$\dfrac{1}{a} > c$，因此$f(x) < 0$的解集为$(c, \dfrac{1}{a})$
                \item 根据题意可知该三角形面积为$S = \dfrac{1}{2}(\dfrac{1}{a} - c) \times c = 8$，故$\dfrac{1}{a} = c + \dfrac{16}{c} \in [8, +\infty)$，因此$a \in (0, \dfrac{1}{8}]$
                \item 根据$f(c) = 0$可知$1+b+ac = 0$，因此不等式变为$m^2-2km \geq 0$，
                以$k$为主元，设$g(k) = m^2 - 2km$，则应有$g(-1) \geq 0, g(1) \geq 0$，解得$m \in (-\infty, -2] \cup [2, +\infty) \cup \{0\}$
            \end{enumerate}
        \end{jieda}
    \end{enumerate}
    \item 根式和的最值与柯西不等式
    \begin{enumerate}
        \item 函数$y = \sqrt{x} + \sqrt{1-x}$的最大值为\blkda{$\sqrt{2}$}
        \begin{jieda}
            \begin{enumerate}
            \item 平方处理 \\
            $y^2 = x + 1 - x + 2\sqrt{x(1-x)} = 1 + 1 + 2\sqrt{x(1-x)} \leq 2$，故$y \leq \sqrt{2}$
            \item 利用平均值不等式链 \\
            根据$\dfrac{a+b}{2} \leq \sqrt{\dfrac{a^2+b^2}{2}}$可知$y = \sqrt{x} + \sqrt{1-x} \leq 2\sqrt{\dfrac{x + 1 - x}{2}} = \sqrt{2}$
            \item 利用柯西不等式 \\
            $y = \sqrt{x} + \sqrt{1-x} = 1 \times \sqrt{x} + 1 \times \sqrt{1-x}$，利用柯西不等式（$(a^2+b^2)\cdot(c^2+d^2) \geq (ac+bd)^2$）有$y \leq \sqrt{(1+1) \cdot (x + 1-x)} = \sqrt{2}$，在$\sqrt{1-x} = \sqrt{x}$时，即$x = \dfrac{1}{2}$时取等（在定义域$[0,1]$内）
            \item 数形结合 \\
            设$a = \sqrt{x}, b = \sqrt{1-x} (a \in [0, 1], b \in [0, 1])$，则两者满足$a^2 + b^2 = 1$，为单位圆方程($a$为横坐标，$b$为纵坐标)。$y = a+b$为一直线方程($a$为横坐标，$b$为纵坐标)，故$y$取最大值的时候为直线和圆相切。当直线与圆相切时，根据点到直线的距离公式有原点到直线的距离$\dfrac{|y|}{\sqrt{1+1}} = 1$，因此$y = \sqrt{2}$
            \item 三角换元 \\
            设$cos\theta = \sqrt{x}, sin\theta = \sqrt{1-x} \left(\theta \in \left[0, \dfrac{\pi}{2}\right]\right)$，故$y = sin \theta + cos\theta = \sqrt{2}\left(sin\left(\theta + \dfrac{\pi}{4}\right)\right) \leq \sqrt{2}$，当$\theta = \dfrac{\pi}{4}$时取得最值.
        \end{enumerate}        
        \end{jieda}
        \item 函数$y = 4\sqrt{x-2} + \sqrt{9-3x}$的最大值为\blkda{$\sqrt{19}$}
        \begin{jieda}
            \begin{enumerate}
                \item 利用柯西不等式 \\
                $y = 4\sqrt{x-2} + \sqrt{9-3x} = 4\sqrt{x-2} + \sqrt{3}\sqrt{3-x}$，利用柯西不等式（$(a^2+b^2)\cdot(c^2+d^2) \geq (ac+bd)^2$）有$y \leq \sqrt{(16+3) \cdot (x-2 + 3-x)} = \sqrt{19}$，在$\sqrt{3}\sqrt{3-x} = 4\sqrt{x-2}$时，即$x = \dfrac{54}{19}$时取等（在定义域$[2,3]$内）
                \item 数形结合 \\
                \ding{172} 设$a = 4\sqrt{x-2}, b = \sqrt{9-3x} (a \in [0, 4], b \in [0, \sqrt{3}])$，则两者满足$3a^2 + 16b^2 = 48$，即$\dfrac{a^2}{16} + \dfrac{b^2}{3} = 1$，为一椭圆方程($a$为横坐标，$b$为纵坐标)。$y = a+b$为一直线方程($a$为横坐标，$b$为纵坐标)，故$y$取最大值的时候为直线和椭圆相切。故此处可以采用隐函数求导的方式得到$(a_0, b_0)$处的切线方程：$\dfrac{2a_0a}{16} + \dfrac{2b_0b}{3} = 1$，故当斜率为$-1$时$3a_0 = 16b_0$，代入椭圆方程得到$b_0 = \dfrac{3}{\sqrt{19}}$，故$a_0 = \dfrac{16}{\sqrt{19}}$，因此$y = a_0 + b_0 = \sqrt{19}$ \\
                \ding{173} 也可以设$c = 4\sqrt{3}\sqrt{x-2}, d = 4\sqrt{9-3x}$，满足$c^2+d^2 = 48, y = \dfrac{c}{\sqrt{3}} + \dfrac{d}{4}$.\\
                $y$取最大值的时候为直线和圆相切。故利用点到直线的距离公式有$\sqrt{48} = \dfrac{y}{\sqrt{\dfrac{1}{3} + \dfrac{1}{16}}}$，\\
                即$y = \sqrt{48} \cdot \sqrt{\dfrac{19}{48}} = \sqrt{19}$
            \end{enumerate}
        \end{jieda}
    \end{enumerate}
    \item *万能$k$法
    \begin{enumerate}
        \item 已知正实数$x, y$满足$xy^2(x+y) = 4$，则$2x+y$的最小值为\blkda{$2\sqrt{2}$}
        \begin{jieda}
            观察可知$xy^2(x+y) = 4$不能够将$x, y$分离，且$2x+y$为一次相加的形式。此时应考虑设$2x+y = k$，代入$xy^2(x+y) = 4$中求$k$的最小值即可。因为$xy^2(x+y) = 4$中$x$为二次，$y$为三次，因此考虑代换$x$，即$x = \dfrac{k-y}{2}$ \\
            $\dfrac{k-y}{2}y^2(\dfrac{k+y}{2}) = 4$，即$(k-y)y^2(k+y) = 16$，进一步整理得$k^2y^2 - y^4 = 16$，参变分离得到$k^2 = \dfrac{16}{y^2} + y^2$，根据对勾函数性质知$k^2$最小值为8，故$k$最小值为$2\sqrt{2}$（$y = 2, x = \sqrt{2}-1$时取最小值.）
        \end{jieda}
        \item 设$x, y$为实数，若$4x^2 + y^2 -xy = 3$，则$2x+y$的最大值为\blkda{$2\sqrt{2}$}
        \begin{jieda}
            设$k = 2x+y$，则$4x^2 + (k-2x)^2 - x(k-2x) - 3 = 0$，整理有$10x^2 - 5kx + k^2 -3 = 0$，方程有解故$\Delta = 25k^2 - 40(k^2-3) = 120-15k^2 \geq 0$，故$k \in [-2\sqrt{2}, 2\sqrt{2}]$，因此$2x+y$的最大值为$2\sqrt{2}$
        \end{jieda}
        \item 已知实数$a, b$满足$a^3 + b^3 + 3ab = 1$，设$a+b$所有可能取值构成的集合为$M$，则\dotfill{(\kh{B})}
        \xx{$M$为单元素集}{$M$为有限集，但不是单元素集}{$M$为无限集，且有下界}{$M$为无限集，且无下界}
        \begin{jieda}
            \begin{enumerate}
                \item 考虑使用万能$k$法 \\
                设$k = a+b$，则$a^3 + b^3 + 3ab = 1$可以化为$a^3 + (k-a)^3 + 3a(k-a) = 1$ \\
                进一步打开括号有$a^3 + k^3 - 3k^2a + 3ka^2 - a^3 + 3ka - 3a^2 = 1$，\\
                合并同类项有$(3k-3)a^2 + (3k-3k^2)a + k^3 - 1 = 0$，\\
                提取公因式后有$(k-1)[3a^2 - 3ka + (k^2+k+1)] = 0$，\\
                故必有$1 \in M$.\\
                当$k \neq 1$时，考虑$3a^2 - 3ka + (k^2+k+1) = 0$，\\
                将其看做关于$k$的二次方程，整理得$k^2 + (1-3a)k + 1 + 3a^2 = 0$ \\
                其判别式为$\Delta = -3(a+1)^2$，故仅当$a = -1$时有解$k = -2$，故$M = \{-2, 1\}$
                \item 考虑利用完全立方公式进行凑配。\\
                $a^3 + b^3 + 3ab = 1$可以化为$(a+b)^3 -3a^2b - 3ab^2 + 3ab = 1$ \\
                进一步提取公因式有$(a+b)^3 - 1 = 3ab(a+b+1)$，\\
                左侧利用立方差公式有$(a+b-1)((a+b)^2+1+(a+b)) = 3ab(a+b-1)$ \\
                由此可以看到必有$1 \in M$.\\
                当$a+b \neq 1$时$(a+b)^2+1+(a+b) = 3ab$ \\
                进一步凑配$a+b$的形式，有$(a+b)^2+1+(a+b) = 3ab + 3a^2 - 3a^2 = 3a(a+b) - 3a^2$，即$(a+b)^2+(a+b)(1-3a)+1+3a^2 = 0$\\
                其判别式为$\Delta = -3(a+1)^2$，故仅当$a = -1$时有解$a+b = -2$，故$M = \{-2, 1\}$
            \end{enumerate}
        \end{jieda}
    \end{enumerate}
    \item *选择主元求取值范围
    \begin{enumerate}
        \item 已知实数$x, y, z$满足$x^2+y^2+z^2+xy+yz+zx = 1$，则下列说法错误的是\dotfill{(\kh{A})}
        \xx{$xyz$的最大值为$\dfrac{\sqrt{6}}{6}$}{$x+y+z$的最大值为$\dfrac{\sqrt{6}}{2}$}{$x$的最大值为$\dfrac{\sqrt{6}}{2}$}{$x+y$的最大值为$\sqrt{2}$}
        \begin{jieda}
            \begin{enumerate}[A.]
                \item 利用拉格朗日乘数法可求最值(\ps 高中范围内知识应该难以求出最值) \\
                设$f(x,y,z) = xyz - \lambda \cdot (x^2+y^2+z^2+xy+yz+zx - 1)$，令偏导数为0： \\
                $\dfrac{\partial f}{\partial x} = yz - \lambda \cdot (2x+y+z) = 0$，$\dfrac{\partial f}{\partial y} = xz - \lambda \cdot (2y+x+z) = 0$，$\dfrac{\partial f}{\partial z} = xy - \lambda \cdot (2z+x+y) = 0$ \\
                $\left \{ \begin{array}{l}
                     yz = \lambda \cdot (2x+y+z)  \\
                     xz = \lambda \cdot (2y+x+z)
                \end{array}\right . \Longrightarrow (y-x)\cdot(z+\lambda) = 0 \Longrightarrow y = z \mbox{或} z = -\lambda$ \\
                $\left \{ \begin{array}{l}
                     yz = \lambda \cdot (2x+y+z)  \\
                     xy = \lambda \cdot (2z+x+y)
                \end{array}\right . \Longrightarrow (x-z)\cdot(y+\lambda) = 0\Longrightarrow x = z \mbox{或} y = -\lambda$
                \begin{dingautolist}{172}
                    \item $x = y = z$ \\
                    此时有$f(x,y,z) = \pm \dfrac{\sqrt{6}}{36}$
                    \item $y = z = -\lambda$ \\
                    $yz = \lambda \cdot (2x+y+z)$即$\lambda^2 = \lambda(2x-2\lambda)$，解得$x = \dfrac{3}{2}\lambda$或$\lambda = 0$ \\
                    若$\lambda = 0$，则$x = \pm 1, xyz = 0$；
                    若$x = \dfrac{3}{2}\lambda$，则代入原条件有$\lambda = \pm \dfrac{2}{3}, xyz = \dfrac{3}{2}\lambda^3 = \pm \dfrac{4}{9}$
                \end{dingautolist}
                综上，$xyz$的最大值为$\dfrac{4}{9}$，$x = y = -\dfrac{2}{3}, z = 1$时可取最值.
                \item $x^2+y^2+z^2+xy+yz+zx = 1$即$(x+y+z)^2 = 1 + xy + yz + zx \leq 1 + (\dfrac{x+y}{2})^2 + (\dfrac{y+z}{2})^2 + (\dfrac{z+x}{2})^2 = \dfrac{3}{2}$，故$x+y+z \leq \dfrac{\sqrt{6}}{2}$（$x = y = z = \dfrac{\sqrt{6}}{6}$时取等.）
                \item 选择$y$为主元，看做关于$y$的方程，即$y^2+(x+z)y+(zx+x^2+z^2-1) = 0$，方程有解则$\Delta = (x+z)^2 - 4(zx+x^2+z^2-1) \geq 0$，整理有$-3x^2-3z^2-2zx+4 \geq 0$，选择$z$为主元，看做关于$z$的不等式，即$3z^2 + 2xz + (3x^2-4) \leq 0$，不等式有解则$\Delta = 4x^2 - 12(3x^2-4) \geq 0$，即$-32x^2 + 48 \geq 0$，故$x \in [-\dfrac{\sqrt{6}}{2}, \dfrac{\sqrt{6}}{2}]$
                \item $x^2+y^2+z^2+xy+yz+zx = 1$即$(x+y)^2+(y+z)^2+(z+x)^2 = 2$，故$x+y \leq \sqrt{2}$（$x = y = \dfrac{\sqrt{2}}{2}, z = 0$时取等号）
            \end{enumerate}
        \end{jieda}
        \item 已知三个互不相同的实数$a, b, c$满足$a+b+c = 1$和$a^2+b^2+c^2 = 3$，则$b$的取值范围是\blkda{$(-1, \dfrac{5}{3})$},已知函数$y = x^3-x^2-x$在$x \in (-\infty, -\dfrac{1}{3}]$和$x \in [1, +\infty)$时$y$随$x$增大而增大，在$x \in (-\dfrac{1}{3}, 1)$时$y$随$x$增大而减小，则$abc$的取值范围是\blkda{$\dfrac{5}{27}$}
        \begin{jieda}
            首先考虑消元：$a+b+c = 1 \Longrightarrow c = 1 - a - b \Longrightarrow c^2 = 1 + a^2 + b^2 - 2a - 2b + 2ab$\\
            $c^2 = 3 - a^2 - b^2$，故$1 + a^2 + b^2 - 2a - 2b + 2ab = 3 - a^2 - b^2$，即$a^2 + b^2 - a - b + ab = 1$ 
            \begin{dingautolist}{172}
                \item 以$a$为主元，看做关于$a$的方程，有$a^2 + (b-1)a + (b^2-b-1) = 0$，方程有解，且由轮换式性质可知重根情况对应着$a, b$相等，不符要求，因此有$\Delta = (b-1)^2 - 4(b^2-b-1) > 0$，即$-3b^2+2b+5 \geq 0$，解得$b \in (-1, \dfrac{5}{3})$，
                \item $a^2 + b^2 - a - b + ab = 1 \Longrightarrow ab = (a+b)^2 - (a+b) \Longrightarrow 1 + ab = (1-c)^2 - (1-c) = c^2 - c \Longrightarrow abc = c^3 - c^2 - c$，
                由轮换式性质可知，$c \in (-1, \dfrac{5}{3})$，故$c = -\dfrac{1}{3}$时$abc$取最大值$\dfrac{5}{27}$
            \end{dingautolist}
            
        \end{jieda}
    \end{enumerate}
    \item *不等式性质与最值函数
    \begin{enumerate}
        \item 定义$max M$表示数集$M$中最大的数，则$max\left\{\dfrac{1}{a}, \dfrac{1}{b}, a^2+b^2\right\}$的最小值为\blkda{$\sqrt[3]{2}$}
        \begin{jieda}
            考虑特例$\dfrac{1}{a} = \dfrac{1}{b} = a^2+b^2$，此时$a = b = \sqrt[3]{\dfrac{1}{2}}$，$max\left\{\dfrac{1}{a}, \dfrac{1}{b}, a^2+b^2\right\} = \sqrt[3]{2}$ 下面反证法证明$max\left\{\dfrac{1}{a}, \dfrac{1}{b}, a^2+b^2\right\}$取得最小值为$\sqrt[3]{2}$. \\
            假设$max\left\{\dfrac{1}{a}, \dfrac{1}{b}, a^2+b^2\right\}$的最小值为$m < \sqrt[3]{2}$
            \begin{dingautolist}{172}
                \item 设$\dfrac{1}{a} = m, \dfrac{1}{b} \leq m, a^2+b^2 \leq m$ \\
                此时$a = \dfrac{1}{m}, b \geq \dfrac{1}{m}$，故$a^2+b^2 \geq \dfrac{2}{m^2} > \sqrt[3]{2} > m$，矛盾.
                \item 设$\dfrac{1}{a} \leq m, \dfrac{1}{b} = m, a^2+b^2 \leq m$\\
                此时$a \geq \dfrac{1}{m}, b = \dfrac{1}{m}$，故$a^2+b^2 \geq \dfrac{2}{m^2} > \sqrt[3]{2} > m$，矛盾.
                \item 设$\dfrac{1}{a} \leq m, \dfrac{1}{b} \leq m, a^2+b^2 = m$ \\
                此时$a \geq \dfrac{1}{m}, b \geq \dfrac{1}{m}$，故$a^2+b^2 \geq \dfrac{2}{m^2} > \sqrt[3]{2} > m$，矛盾.
            \end{dingautolist}
        \end{jieda}
        \item 定义$max M$表示数集$M$中最大的数.设$0<a<b<c<1$，已知$b\geq 2a$或$a+b\leq 1$，则$max\left\{b-a, c-b, 1-c\right\}$的最小值为\blkda{$\dfrac{1}{5}$}
        \begin{jieda}
            设$A = a, B = b-a, C = c - b, D = 1 - c$，则$A+B+C+D = 1$且有$A \leq B$或$2A+B \leq 1$，目标值为$max\left\{B, C, D\right\}$

            考虑特例$B = C = D$，则此时有$A+3B = 1$，故$A = 1-3B$
            
            因此$A \leq B$或$2A+B \leq 1$即$1-3B \leq B$或$2-5B \leq 1$，即$B \geq \dfrac{1}{4}$或$B \geq \dfrac{1}{5}$，故有特例$A = \dfrac{2}{5}, B = C = D = \dfrac{1}{5}$
            
            下面反证法证明此特例是使得$max\left\{B, C, D\right\}$取得最小值的例子.
    
            
            假设$max\left\{B, C, D\right\}$可取到最小值$m < \dfrac{1}{5}$，不失一般性可以假设$B = m < \dfrac{1}{5}, C \leq m, D \leq m$，此时\ding{172}$A \geq 1 - 3m > \dfrac{2}{5}$，故$A > B$，\ding{173}$2A + B = 2 - 5m > 1$，故$2A+B > 1$，矛盾.
            因此$max\left\{b-a, c-b, 1-c\right\}$的最小值为$\dfrac{1}{5}$
        \end{jieda}
        \item 定义$max M$表示数集$M$中最大的数，设$0 < a < b < c < 2$，已知$b \geq 3a$或$a+2b \leq 3$，则$max\left\{a-b, c-b, 2-c\right\}$的最小值为\blkda{$\dfrac{3}{7}$}
        \begin{jieda}
            设$A = a, B = b-a, C = c - b, D = 2 - c$，则$A+B+C+D = 2$且有$2A \leq B$或$3A+2B \leq 3$ \\
            考虑特例$B = C = D$，则此时有$A+3B = 2$，故$A = 2-3B$，$2A \leq B$或$3A+2B \leq 3$即$2-3B \leq B$或$6-7B \leq 3$，即$B \geq \dfrac{1}{2}$或$B \geq \dfrac{3}{7}$，故有特例$A = \dfrac{5}{7}, B = C = D = \dfrac{3}{7}$，下面反证法证明此特例是使得$max\left\{B, C, D\right\}$取得最小值的例子. \\
            假设$max\left\{B, C, D\right\}$可取到最小值$m < \dfrac{3}{7}$，不失一般性可以假设$B = m < \dfrac{3}{7}, C \leq m, D \leq m$，此时\ding{172}$A \geq 2 - 3m > \dfrac{5}{7}$，故$2A > B$，\ding{173}$3A + 2B = 6 - 7m > 3$，故$3A+2B > 3$，矛盾.
            因此$max\left\{b-a, c-b, 1-c\right\}$的最小值为$\dfrac{3}{7}$
        \end{jieda}
    \end{enumerate}
\end{enumerate}